\begin{PSbox}[unbreakable]{obj}{اهداف یادگیری}

پس از تکمیل این ماژول، دانشجویان قادر خواهند بود:

\begin{enumerate}
\def\labelenumi{\arabic{enumi}.}
\tightlist
\item
  تشخیص دهند که کدام توزیع احتمال یک سناریوی مهندسی معین را مدل می‌کند
\item
  \lr{$\frac{\mathrm{PMF}}{\mathrm{PDF}}$}، \lr{\mbox{CDF}}، میانگین و واریانس را برای تمام توزیع‌های اصلی بیان کنند
\item
  تفسیر فیزیکی و زمینه مهندسی هر توزیع را توضیح دهند
\item
  تعیین کنند چه زمانی از هر توزیع استفاده شود (و چه زمانی استفاده نشود)
\item
  توزیع‌ها را برای شبیه‌سازی در \lr{\mbox{MATLAB}} پیاده‌سازی کنند
\item
  الگوهای توزیع را از مشخصه‌های داده تشخیص دهند
\end{enumerate}

\end{PSbox}

\PSsection{3.1}{مقدمه: انتخاب توزیع مناسب}

\PSsubsection{چالش مهندسی}

هنگام مدل‌سازی یک پدیده تصادفی، پرسش حیاتی این است:

\begin{PSquote}

«کدام توزیع احتمال این کمیت نامعین را بهتر توصیف می‌کند؟»

\end{PSquote}

انتخاب توزیع نادرست به پیش‌بینی‌های نادرست، طراحی‌های نامناسب و خرابی سیستم‌ها منجر می‌شود.

\PSsubsection{چگونه انتخاب کنیم}

توزیع باید با \textbf{سازوکار فیزیکی} تولید تصادفی‌بودن همخوانی داشته باشد:

\begin{itemize}
\tightlist
\item
  شمارش موفقیت‌ها \lr{\mbox{\ensuremath{\to}}} دوجمله‌ای
\item
  انتظار برای رویدادها \lr{\mbox{\ensuremath{\to}}} نمایی/هندسی
\item
  جمع بسیاری از اثرهای کوچک \lr{\mbox{\ensuremath{\to}}} گاوسی
\item
  پوش سیگنال در محوشدگی \lr{\mbox{\ensuremath{\to}}} رایلی
\end{itemize}

این ماژول ابزارهای لازم برای انجام درست این انتخاب را فراهم می‌کند.

\PSsection{3.2}{توزیع برنولی}

\PSsubsection{تفسیر فیزیکی}

یک \textbf{آزمایش منفرد} با دو خروجی ممکن را مدل می‌کند: «موفقیت» (1) یا «شکست» (0).

\PSsubsection{زمینه مهندسی}

هر خروجی تصادفی دودویی:

\begin{itemize}
\tightlist
\item
  یک بیت ارسال‌شده: خطا یا درست
\item
  یک قطعه آزمایش‌شده: قبول یا رد
\item
  یک بسته ارسال‌شده: تحویل‌شده یا مفقود
\item
  یک تلاش آشکارسازی: هدف حاضر یا غایب
\end{itemize}

\begin{PSbox}[unbreakable]{def}{تعریف ریاضی}

\lr{$X \sim \mathrm{Bernoulli}(p)$}، که در آن \lr{$p$} = احتمال موفقیت.

\end{PSbox}

\PSsubsection{\lr{\mbox{PMF}}}

\PSdisplay{p_X(x) = \begin{cases} p & x = 1 \\ 1-p = q & x = 0 \end{cases}}

شکل فشرده: \lr{$p_{X}(x) = p^{x}(1-p)^{1-x}$} برای \lr{$x \in \{0,\, 1\}$}

\PSsubsection{\lr{\mbox{CDF}}}

\PSdisplay{F_X(x) = \begin{cases} 0 & x < 0 \\ 1-p & 0 \leq x < 1 \\ 1 & x \geq 1 \end{cases}}

\PSsubsection{میانگین}

\PSdisplay{E[X] = p}

\PSsubsection{واریانس}

\PSdisplay{\PSmtext{Var}(X) = p(1-p)}

بیشترین واریانس در \lr{$p = 0.5$} است (بیشترین عدم‌قطعیت).

\PSsubsection{زمان استفاده}

\lr{\mbox{\PSyes{}}} آزمایش دودویی منفرد با احتمال ثابت\PSbr{}\lr{\mbox{\PSyes{}}} بلوک سازنده برای توزیع‌های پیچیده‌تر\PSbr{}\lr{\mbox{\PSyes{}}} مدل‌سازی حالت‌های روشن/خاموش، خروجی‌های قبول/رد

\PSsubsection{زمان عدم استفاده}

\lr{\mbox{\PSno{}}} آزمایش‌های متعدد (از دوجمله‌ای استفاده کنید)\PSbr{}\lr{\mbox{\PSno{}}} خروجی دودویی نیست\PSbr{}\lr{\mbox{\PSno{}}} احتمال میان آزمایش‌ها تغییر می‌کند

\PSsubsection{پیاده‌سازی در \lr{\mbox{MATLAB}}}

\tcbset{PScodemode/.style={unbreakable}}

\begin{Shaded}
\begin{Highlighting}[]
\CommentTok{\% Generate Bernoulli random variables}
\VariableTok{p} \OperatorTok{=} \FloatTok{0.1}\OperatorTok{;}          \CommentTok{\% Bit error probability}
\VariableTok{N} \OperatorTok{=} \FloatTok{10000}\OperatorTok{;}
\VariableTok{X} \OperatorTok{=} \VariableTok{binornd}\NormalTok{(}\FloatTok{1}\OperatorTok{,} \VariableTok{p}\OperatorTok{,} \FloatTok{1}\OperatorTok{,} \VariableTok{N}\NormalTok{)}\OperatorTok{;}  \CommentTok{\% or: X = rand(1,N) \textless{} p;}

\CommentTok{\% Verify}
\VariableTok{fprintf}\NormalTok{(}\SpecialStringTok{\textquotesingle{}Mean: \%.4f (theory: \%.4f)\textbackslash{}n\textquotesingle{}}\OperatorTok{,} \VariableTok{mean}\NormalTok{(}\VariableTok{X}\NormalTok{)}\OperatorTok{,} \VariableTok{p}\NormalTok{)}\OperatorTok{;}
\VariableTok{fprintf}\NormalTok{(}\SpecialStringTok{\textquotesingle{}Var:  \%.4f (theory: \%.4f)\textbackslash{}n\textquotesingle{}}\OperatorTok{,} \VariableTok{var}\NormalTok{(}\VariableTok{X}\NormalTok{)}\OperatorTok{,} \VariableTok{p}\OperatorTok{*}\NormalTok{(}\FloatTok{1}\OperatorTok{{-}}\VariableTok{p}\NormalTok{))}\OperatorTok{;}
\end{Highlighting}
\end{Shaded}

\PSsection{3.3}{توزیع دوجمله‌ای}

\PSsubsection{تفسیر فیزیکی}

\textbf{تعداد موفقیت‌ها} را در \lr{$n$} آزمایش \textbf{مستقل} که هر یک همان احتمال موفقیت \lr{$p$} را دارند مدل می‌کند.

\PSsubsection{زمینه مهندسی}

\begin{itemize}
\tightlist
\item
  تعداد خطاهای بیت در \lr{$n$} بیت ارسال‌شده
\item
  تعداد قطعات معیوب در یک دسته \lr{\mbox{n}}‌تایی
\item
  تعداد انتقال‌های موفق بسته از میان \lr{$n$} تلاش
\item
  تعداد حسگرهایی که یک سیگنال را آشکار می‌کنند (از میان \lr{$n$} حسگر)
\end{itemize}

\begin{PSbox}[unbreakable]{def}{تعریف ریاضی}

\lr{$X \sim \mathrm{Binomial}(n,\, p)$}

\end{PSbox}

\PSsubsection{\lr{\mbox{PMF}}}

\PSdisplay{p_X(k) = \binom{n}{k} p^k (1-p)^{n-k}, \quad k = 0, 1, 2, \ldots, n}

که در آن \lr{$C(n,\,k) = \frac{n!}{k!(n-k)!}$}

\PSsubsection{\lr{\mbox{CDF}}}

\PSdisplay{F_X(k) = \sum_{i=0}^{\lfloor k \rfloor} \binom{n}{i} p^i (1-p)^{n-i}}

فرم بسته ندارد \lr{—} از جدول‌ها یا \lr{\mbox{MATLAB}} استفاده کنید.

\PSsubsection{میانگین}

\PSdisplay{E[X] = np}

\PSsubsection{واریانس}

\PSdisplay{\PSmtext{Var}(X) = np(1-p)}

\PSsubsection{زمان استفاده}

\lr{\mbox{\PSyes{}}} تعداد ثابت آزمایش‌ها \lr{$n$}\PSbr{}\lr{\mbox{\PSyes{}}} هر آزمایش مستقل است\PSbr{}\lr{\mbox{\PSyes{}}} هر آزمایش همان احتمال \lr{$p$} را دارد\PSbr{}\lr{\mbox{\PSyes{}}} شمارش تعداد «موفقیت‌ها»

\PSsubsection{زمان عدم استفاده}

\lr{\mbox{\PSno{}}} آزمایش‌ها مستقل نیستند (خرابی‌های همبسته)\PSbr{}\lr{\mbox{\PSno{}}} احتمال میان آزمایش‌ها تغییر می‌کند\PSbr{}\lr{\mbox{\PSno{} $n$}} ثابت نیست (از دوجمله‌ای منفی یا پواسون استفاده کنید)\PSbr{}\lr{\mbox{\PSno{} $n$}} بسیار بزرگ با \lr{$p$} کوچک (از تقریب پواسون استفاده کنید)

\PSsubsection{پیاده‌سازی در \lr{\mbox{MATLAB}}}

\tcbset{PScodemode/.style={unbreakable}}

\begin{Shaded}
\begin{Highlighting}[]
\CommentTok{\%\% Binomial Distribution: Bit Errors in a Frame}
\VariableTok{n} \OperatorTok{=} \FloatTok{100}\OperatorTok{;}         \CommentTok{\% bits per frame}
\VariableTok{p} \OperatorTok{=} \FloatTok{0.05}\OperatorTok{;}        \CommentTok{\% bit error probability}
\VariableTok{k} \OperatorTok{=} \FloatTok{0}\OperatorTok{:}\FloatTok{20}\OperatorTok{;}        \CommentTok{\% range of interest}

\CommentTok{\% PMF and CDF}
\VariableTok{pmf} \OperatorTok{=} \VariableTok{binopdf}\NormalTok{(}\VariableTok{k}\OperatorTok{,} \VariableTok{n}\OperatorTok{,} \VariableTok{p}\NormalTok{)}\OperatorTok{;}
\VariableTok{cdf\_vals} \OperatorTok{=} \VariableTok{binocdf}\NormalTok{(}\VariableTok{k}\OperatorTok{,} \VariableTok{n}\OperatorTok{,} \VariableTok{p}\NormalTok{)}\OperatorTok{;}

\CommentTok{\% Plot PMF}
\VariableTok{figure}\OperatorTok{;}
\VariableTok{stem}\NormalTok{(}\VariableTok{k}\OperatorTok{,} \VariableTok{pmf}\OperatorTok{,} \SpecialStringTok{\textquotesingle{}filled\textquotesingle{}}\OperatorTok{,} \SpecialStringTok{\textquotesingle{}LineWidth\textquotesingle{}}\OperatorTok{,} \FloatTok{1.5}\NormalTok{)}\OperatorTok{;}
\VariableTok{xlabel}\NormalTok{(}\SpecialStringTok{\textquotesingle{}Number of Errors (k)\textquotesingle{}}\NormalTok{)}\OperatorTok{;}
\VariableTok{ylabel}\NormalTok{(}\SpecialStringTok{\textquotesingle{}P(X = k)\textquotesingle{}}\NormalTok{)}\OperatorTok{;}
\VariableTok{title}\NormalTok{(}\VariableTok{sprintf}\NormalTok{(}\SpecialStringTok{\textquotesingle{}Binomial PMF: n=\%d, p=\%.2f\textquotesingle{}}\OperatorTok{,} \VariableTok{n}\OperatorTok{,} \VariableTok{p}\NormalTok{))}\OperatorTok{;}
\VariableTok{grid} \VariableTok{on}\OperatorTok{;}

\CommentTok{\% Simulation verification}
\VariableTok{N\_sim} \OperatorTok{=} \FloatTok{100000}\OperatorTok{;}
\VariableTok{X\_sim} \OperatorTok{=} \VariableTok{binornd}\NormalTok{(}\VariableTok{n}\OperatorTok{,} \VariableTok{p}\OperatorTok{,} \FloatTok{1}\OperatorTok{,} \VariableTok{N\_sim}\NormalTok{)}\OperatorTok{;}
\VariableTok{fprintf}\NormalTok{(}\SpecialStringTok{\textquotesingle{}E[X]: Theory = \%.2f, Simulation = \%.2f\textbackslash{}n\textquotesingle{}}\OperatorTok{,} \VariableTok{n}\OperatorTok{*}\VariableTok{p}\OperatorTok{,} \VariableTok{mean}\NormalTok{(}\VariableTok{X\_sim}\NormalTok{))}\OperatorTok{;}
\VariableTok{fprintf}\NormalTok{(}\SpecialStringTok{\textquotesingle{}Var(X): Theory = \%.2f, Simulation = \%.2f\textbackslash{}n\textquotesingle{}}\OperatorTok{,} \VariableTok{n}\OperatorTok{*}\VariableTok{p}\OperatorTok{*}\NormalTok{(}\FloatTok{1}\OperatorTok{{-}}\VariableTok{p}\NormalTok{)}\OperatorTok{,} \VariableTok{var}\NormalTok{(}\VariableTok{X\_sim}\NormalTok{))}\OperatorTok{;}
\end{Highlighting}
\end{Shaded}

\begin{PSbox}[unbreakable]{ex}{مثال شبیه‌سازی: نرخ خطای بسته}

\tcbset{PScodemode/.style={unbreakable}}

\begin{Shaded}
\begin{Highlighting}[]
\CommentTok{\%\% Engineering Scenario: How many packets out of 50 have errors?}
\CommentTok{\% Each packet has 1000 bits, BER = 1e{-}4}
\CommentTok{\% A packet has an error if ANY bit is in error.}

\VariableTok{BER} \OperatorTok{=} \FloatTok{1e{-}4}\OperatorTok{;}
\VariableTok{bits\_per\_packet} \OperatorTok{=} \FloatTok{1000}\OperatorTok{;}
\VariableTok{n\_packets} \OperatorTok{=} \FloatTok{50}\OperatorTok{;}

\CommentTok{\% P(packet error) = 1 {-} P(all bits correct)}
\VariableTok{p\_packet\_error} \OperatorTok{=} \FloatTok{1} \OperatorTok{{-}}\NormalTok{ (}\FloatTok{1} \OperatorTok{{-}} \VariableTok{BER}\NormalTok{)}\OperatorTok{\^{}}\VariableTok{bits\_per\_packet}\OperatorTok{;}
\VariableTok{fprintf}\NormalTok{(}\SpecialStringTok{\textquotesingle{}P(packet error) = \%.4f\textbackslash{}n\textquotesingle{}}\OperatorTok{,} \VariableTok{p\_packet\_error}\NormalTok{)}\OperatorTok{;}

\CommentTok{\% Number of errored packets \textasciitilde{} Binomial(50, p\_packet\_error)}
\VariableTok{k} \OperatorTok{=} \FloatTok{0}\OperatorTok{:}\FloatTok{10}\OperatorTok{;}
\VariableTok{pmf\_packets} \OperatorTok{=} \VariableTok{binopdf}\NormalTok{(}\VariableTok{k}\OperatorTok{,} \VariableTok{n\_packets}\OperatorTok{,} \VariableTok{p\_packet\_error}\NormalTok{)}\OperatorTok{;}
\VariableTok{fprintf}\NormalTok{(}\SpecialStringTok{\textquotesingle{}P(0 packet errors) = \%.4f\textbackslash{}n\textquotesingle{}}\OperatorTok{,} \VariableTok{pmf\_packets}\NormalTok{(}\FloatTok{1}\NormalTok{))}\OperatorTok{;}
\VariableTok{fprintf}\NormalTok{(}\SpecialStringTok{\textquotesingle{}P(≤ 2 packet errors) = \%.4f\textbackslash{}n\textquotesingle{}}\OperatorTok{,} \VariableTok{sum}\NormalTok{(}\VariableTok{pmf\_packets}\NormalTok{(}\FloatTok{1}\OperatorTok{:}\FloatTok{3}\NormalTok{)))}\OperatorTok{;}
\end{Highlighting}
\end{Shaded}

\end{PSbox}

\PSsection{3.4}{توزیع هندسی}

\PSsubsection{تفسیر فیزیکی}

\textbf{تعداد آزمایش‌ها تا نخستین موفقیت} را مدل می‌کند (یا به‌طور معادل، تعداد شکست‌ها پیش از نخستین موفقیت).

\PSsubsection{زمینه مهندسی}

\begin{itemize}
\tightlist
\item
  تعداد انتقال‌ها تا تحویل موفق یک بسته
\item
  تعداد قطعات آزمایش‌شده تا یافتن یک قطعه معیوب
\item
  تعداد تلاش‌های ورود تا احراز هویت موفق
\item
  تعداد شکاف‌های زمانی تا در دسترس قرار گرفتن کانال
\end{itemize}

\begin{PSbox}[unbreakable]{def}{تعریف ریاضی}

\lr{$X \sim \mathrm{Geometric}(p)$}

قرارداد: \lr{$X$} = شماره آزمایشی که نخستین موفقیت در آن رخ می‌دهد \lr{$(X \in \{1,\, 2,\, 3,\, \ldots \})$}.

\end{PSbox}

\PSsubsection{\lr{\mbox{PMF}}}

\PSdisplay{p_X(k) = (1-p)^{k-1} p, \quad k = 1, 2, 3, \ldots}

(\lr{$k-1$} شکست، سپس یک موفقیت)

\PSsubsection{\lr{\mbox{CDF}}}

\PSdisplay{F_X(k) = 1 - (1-p)^k, \quad k = 1, 2, 3, \ldots}

\PSsubsection{میانگین}

\PSdisplay{E[X] = \frac{1}{p}}

\PSsubsection{واریانس}

\PSdisplay{\PSmtext{Var}(X) = \frac{1-p}{p^2}}

\PSsubsection{خاصیت بی‌حافظه}

توزیع هندسی تنها توزیع گسسته با خاصیت بی‌حافظه است:

\PSdisplay{P(X > m + n \mid X > m) = P(X > n)}

\textbf{معنای مهندسی:} اگر یک سیستم در \lr{$m$} تلاش موفق نشده باشد، احتمال نیاز به \lr{$n$} تلاش بیشتر همانند شروع از صفر است. «شکست‌های گذشته بر احتمال آینده اثر نمی‌گذارند.»

\PSsubsection{زمان استفاده}

\lr{\mbox{\PSyes{}}} آزمایش‌های مستقل مکرر تا نخستین موفقیت\PSbr{}\lr{\mbox{\PSyes{}}} هر آزمایش همان احتمال را دارد\PSbr{}\lr{\mbox{\PSyes{}}} فرض بی‌حافظه مناسب است (مثلاً ارسال‌های مجدد مستقل)

\PSsubsection{زمان عدم استفاده}

\lr{\mbox{\PSno{}}} احتمال موفقیت در طول زمان تغییر می‌کند (پیری، یادگیری)\PSbr{}\lr{\mbox{\PSno{}}} آزمایش‌ها مستقل نیستند\PSbr{}\lr{\mbox{\PSno{}}} به‌دنبال \lr{\mbox{r}}-امین موفقیت هستید (از دوجمله‌ای منفی استفاده کنید)

\PSsubsection{پیاده‌سازی در \lr{\mbox{MATLAB}}}

\tcbset{PScodemode/.style={unbreakable}}

\begin{Shaded}
\begin{Highlighting}[]
\CommentTok{\%\% Geometric Distribution: Retransmissions Until Success}
\VariableTok{p\_success} \OperatorTok{=} \FloatTok{0.7}\OperatorTok{;}   \CommentTok{\% P(successful transmission)}
\VariableTok{k} \OperatorTok{=} \FloatTok{1}\OperatorTok{:}\FloatTok{15}\OperatorTok{;}

\CommentTok{\% PMF}
\VariableTok{pmf\_geom} \OperatorTok{=} \VariableTok{geopdf}\NormalTok{(}\VariableTok{k}\OperatorTok{{-}}\FloatTok{1}\OperatorTok{,} \VariableTok{p\_success}\NormalTok{)}\OperatorTok{;}  \CommentTok{\% MATLAB counts failures (k{-}1)}

\CommentTok{\% Simulation}
\VariableTok{N} \OperatorTok{=} \FloatTok{100000}\OperatorTok{;}
\VariableTok{X\_sim} \OperatorTok{=} \VariableTok{geornd}\NormalTok{(}\VariableTok{p\_success}\OperatorTok{,} \FloatTok{1}\OperatorTok{,} \VariableTok{N}\NormalTok{) }\OperatorTok{+} \FloatTok{1}\OperatorTok{;}  \CommentTok{\% +1 to count trial number}

\VariableTok{figure}\OperatorTok{;}
\VariableTok{subplot}\NormalTok{(}\FloatTok{1}\OperatorTok{,}\FloatTok{2}\OperatorTok{,}\FloatTok{1}\NormalTok{)}\OperatorTok{;}
\VariableTok{stem}\NormalTok{(}\VariableTok{k}\OperatorTok{,} \VariableTok{pmf\_geom}\OperatorTok{,} \SpecialStringTok{\textquotesingle{}filled\textquotesingle{}}\NormalTok{)}\OperatorTok{;}
\VariableTok{xlabel}\NormalTok{(}\SpecialStringTok{\textquotesingle{}Trial Number k\textquotesingle{}}\NormalTok{)}\OperatorTok{;} \VariableTok{ylabel}\NormalTok{(}\SpecialStringTok{\textquotesingle{}P(X = k)\textquotesingle{}}\NormalTok{)}\OperatorTok{;}
\VariableTok{title}\NormalTok{(}\VariableTok{sprintf}\NormalTok{(}\SpecialStringTok{\textquotesingle{}Geometric PMF (p = \%.1f)\textquotesingle{}}\OperatorTok{,} \VariableTok{p\_success}\NormalTok{))}\OperatorTok{;}
\VariableTok{grid} \VariableTok{on}\OperatorTok{;}

\VariableTok{subplot}\NormalTok{(}\FloatTok{1}\OperatorTok{,}\FloatTok{2}\OperatorTok{,}\FloatTok{2}\NormalTok{)}\OperatorTok{;}
\VariableTok{histogram}\NormalTok{(}\VariableTok{X\_sim}\OperatorTok{,} \SpecialStringTok{\textquotesingle{}Normalization\textquotesingle{}}\OperatorTok{,} \SpecialStringTok{\textquotesingle{}probability\textquotesingle{}}\OperatorTok{,} \SpecialStringTok{\textquotesingle{}BinMethod\textquotesingle{}}\OperatorTok{,} \SpecialStringTok{\textquotesingle{}integers\textquotesingle{}}\NormalTok{)}\OperatorTok{;}
\VariableTok{xlabel}\NormalTok{(}\SpecialStringTok{\textquotesingle{}Trial Number k\textquotesingle{}}\NormalTok{)}\OperatorTok{;} \VariableTok{ylabel}\NormalTok{(}\SpecialStringTok{\textquotesingle{}Relative Frequency\textquotesingle{}}\NormalTok{)}\OperatorTok{;}
\VariableTok{title}\NormalTok{(}\SpecialStringTok{\textquotesingle{}Simulation Histogram\textquotesingle{}}\NormalTok{)}\OperatorTok{;}
\VariableTok{grid} \VariableTok{on}\OperatorTok{;}

\VariableTok{fprintf}\NormalTok{(}\SpecialStringTok{\textquotesingle{}E[X]: Theory = \%.2f, Simulation = \%.2f\textbackslash{}n\textquotesingle{}}\OperatorTok{,} \FloatTok{1}\OperatorTok{/}\VariableTok{p\_success}\OperatorTok{,} \VariableTok{mean}\NormalTok{(}\VariableTok{X\_sim}\NormalTok{))}\OperatorTok{;}
\end{Highlighting}
\end{Shaded}

\PSsection{3.5}{توزیع دوجمله‌ای منفی}

\PSsubsection{تفسیر فیزیکی}

\textbf{تعداد آزمایش‌ها تا \lr{\mbox{r}}-امین موفقیت} را مدل می‌کند (تعمیم هندسی).

\PSsubsection{زمینه مهندسی}

\begin{itemize}
\tightlist
\item
  تعداد انتقال‌ها تا تحویل موفق \lr{$r$} بسته
\item
  تعداد قطعات بازرسی‌شده تا یافتن \lr{$r$} قطعه معیوب
\item
  تعداد شکاف‌های زمانی تا در دسترس قرار گرفتن \lr{$r$} کانال
\end{itemize}

\begin{PSbox}[unbreakable]{def}{تعریف ریاضی}

\lr{$X \sim \mathrm{NegBin}(r,\, p)$} = تعداد آزمایش‌ها تا \lr{\mbox{r}}-امین موفقیت.

\end{PSbox}

\PSsubsection{\lr{\mbox{PMF}}}

\PSdisplay{p_X(k) = \binom{k-1}{r-1} p^r (1-p)^{k-r}, \quad k = r, r+1, r+2, \ldots}

(انتخاب این‌که کدام \lr{$r-1$} آزمایش از میان \lr{$k-1$} آزمایش نخست موفق بوده‌اند؛ آزمایش \lr{\mbox{k}}-ام همان \lr{\mbox{r}}-امین موفقیت است.)

\PSsubsection{\lr{\mbox{CDF}}}

فرم بسته ساده‌ای ندارد. با جمع‌بندی یا \lr{\mbox{MATLAB}} محاسبه می‌شود.

\PSsubsection{میانگین}

\PSdisplay{E[X] = \frac{r}{p}}

\PSsubsection{واریانس}

\PSdisplay{\PSmtext{Var}(X) = \frac{r(1-p)}{p^2}}

\PSsubsection{رابطه با هندسی}

وقتی \lr{$r = 1$}، توزیع دوجمله‌ای منفی به توزیع هندسی کاهش می‌یابد.

\PSsubsection{زمان استفاده}

\lr{\mbox{\PSyes{}}} آزمایش‌های مستقل مکرر تا \lr{\mbox{r}}-امین موفقیت\PSbr{}\lr{\mbox{\PSyes{}}} احتمال ثابت در هر آزمایش\PSbr{}\lr{\mbox{\PSyes{}}} تعمیم هندسی به چندین موفقیت موردنیاز

\PSsubsection{زمان عدم استفاده}

\lr{\mbox{\PSno{}}} فقط نخستین موفقیت لازم است (از هندسی استفاده کنید \lr{—} ساده‌تر)\PSbr{}\lr{\mbox{\PSno{}}} آزمایش‌ها مستقل نیستند\PSbr{}\lr{\mbox{\PSno{}}} تعداد آزمایش‌ها ثابت است (از دوجمله‌ای استفاده کنید)

\PSsubsection{پیاده‌سازی در \lr{\mbox{MATLAB}}}

\tcbset{PScodemode/.style={unbreakable}}

\begin{Shaded}
\begin{Highlighting}[]
\CommentTok{\%\% Negative Binomial: Transmissions Until r Successes}
\VariableTok{r} \OperatorTok{=} \FloatTok{5}\OperatorTok{;}             \CommentTok{\% Need 5 successful deliveries}
\VariableTok{p\_success} \OperatorTok{=} \FloatTok{0.8}\OperatorTok{;}   \CommentTok{\% P(success per trial)}

\CommentTok{\% MATLAB\textquotesingle{}s nbinpdf counts FAILURES before r{-}th success}
\CommentTok{\% X\_matlab = number of failures, so total trials = X\_matlab + r}
\VariableTok{k\_failures} \OperatorTok{=} \FloatTok{0}\OperatorTok{:}\FloatTok{20}\OperatorTok{;}

\VariableTok{pmf\_nb} \OperatorTok{=} \VariableTok{nbinpdf}\NormalTok{(}\VariableTok{k\_failures}\OperatorTok{,} \VariableTok{r}\OperatorTok{,} \VariableTok{p\_success}\NormalTok{)}\OperatorTok{;}

\VariableTok{figure}\OperatorTok{;}
\VariableTok{stem}\NormalTok{(}\VariableTok{k\_failures} \OperatorTok{+} \VariableTok{r}\OperatorTok{,} \VariableTok{pmf\_nb}\OperatorTok{,} \SpecialStringTok{\textquotesingle{}filled\textquotesingle{}}\NormalTok{)}\OperatorTok{;}
\VariableTok{xlabel}\NormalTok{(}\SpecialStringTok{\textquotesingle{}Total Trials Until 5th Success\textquotesingle{}}\NormalTok{)}\OperatorTok{;}
\VariableTok{ylabel}\NormalTok{(}\SpecialStringTok{\textquotesingle{}Probability\textquotesingle{}}\NormalTok{)}\OperatorTok{;}
\VariableTok{title}\NormalTok{(}\VariableTok{sprintf}\NormalTok{(}\SpecialStringTok{\textquotesingle{}Negative Binomial: r=\%d, p=\%.1f\textquotesingle{}}\OperatorTok{,} \VariableTok{r}\OperatorTok{,} \VariableTok{p\_success}\NormalTok{))}\OperatorTok{;}
\VariableTok{grid} \VariableTok{on}\OperatorTok{;}

\CommentTok{\% Simulation}
\VariableTok{N} \OperatorTok{=} \FloatTok{100000}\OperatorTok{;}
\VariableTok{X\_sim} \OperatorTok{=} \VariableTok{nbinrnd}\NormalTok{(}\VariableTok{r}\OperatorTok{,} \VariableTok{p\_success}\OperatorTok{,} \FloatTok{1}\OperatorTok{,} \VariableTok{N}\NormalTok{) }\OperatorTok{+} \VariableTok{r}\OperatorTok{;}  \CommentTok{\% total trials}
\VariableTok{fprintf}\NormalTok{(}\SpecialStringTok{\textquotesingle{}E[trials]: Theory = \%.2f, Simulation = \%.2f\textbackslash{}n\textquotesingle{}}\OperatorTok{,} \VariableTok{r}\OperatorTok{/}\VariableTok{p\_success}\OperatorTok{,} \VariableTok{mean}\NormalTok{(}\VariableTok{X\_sim}\NormalTok{))}\OperatorTok{;}
\end{Highlighting}
\end{Shaded}

\PSsection{3.6}{توزیع پواسون}

\PSsubsection{تفسیر فیزیکی}

\textbf{تعداد رویدادهای رخ‌داده در یک بازه ثابت} (زمان، فضا، مساحت) را مدل می‌کند، هنگامی که رویدادها به‌طور مستقل و با نرخ میانگین ثابت رخ می‌دهند.

\PSsubsection{زمینه مهندسی}

\begin{itemize}
\tightlist
\item
  تعداد فوتون‌های برخوردکننده با یک آشکارساز در \lr{$1 \mu s$}
\item
  تعداد بسته‌های واردشده به یک مسیریاب در 1 ثانیه
\item
  تعداد نقص‌ها در هر \lr{$cm^{2}$} روی یک ویفر سیلیکونی
\item
  تعداد رویدادهای پرتو کیهانی در یک حسگر در هر ساعت
\item
  تعداد جهش‌های تداخل در یک پنجره اندازه‌گیری
\end{itemize}

\begin{PSbox}[unbreakable]{def}{تعریف ریاضی}

\lr{$X \sim \mathrm{Poisson}(\lambda)$}، که در آن \lr{$\lambda$} = میانگین تعداد رویدادها در هر بازه.

\end{PSbox}

\PSsubsection{\lr{\mbox{PMF}}}

\PSdisplay{p_X(k) = \frac{\lambda^k e^{-\lambda}}{k!}, \quad k = 0, 1, 2, \ldots}

\PSsubsection{\lr{\mbox{CDF}}}

\PSdisplay{F_X(k) = e^{-\lambda} \sum_{i=0}^{\lfloor k \rfloor} \frac{\lambda^i}{i!}}

\PSsubsection{میانگین}

\PSdisplay{E[X] = \lambda}

\PSsubsection{واریانس}

\PSdisplay{\PSmtext{Var}(X) = \lambda}

\textbf{ویژگی منحصربه‌فرد:} میانگین برابر واریانس است! این یک بررسی سریع است \lr{—} اگر میانگین نمونه \lr{\mbox{\ensuremath{\approx}}} واریانس نمونه باشد، داده ممکن است پواسون باشد.

\PSsubsection{پواسون به‌عنوان تقریب دوجمله‌ای}

وقتی \lr{$n$} بزرگ و \lr{$p$} کوچک باشد، با \lr{\mbox{$\lambda$ = np}}:

\PSdisplay{\PSmtext{Binomial}(n, p) \approx \PSmtext{Poisson}(\lambda = np)}

قاعده سرانگشتی: \lr{$n \ge 20$} و \lr{$p \le 0.05$}.

\PSsubsection{زمان استفاده}

\lr{\mbox{\PSyes{}}} شمارش رویدادها در یک بازه ثابت\PSbr{}\lr{\mbox{\PSyes{}}} رویدادها مستقل رخ می‌دهند\PSbr{}\lr{\mbox{\PSyes{}}} رویدادها با نرخ میانگین ثابت رخ می‌دهند\PSbr{}\lr{\mbox{\PSyes{}}} رویدادهای نادر با فرصت‌های فراوان (\lr{$n$} بزرگ، \lr{$p$} کوچک)\PSbr{}\lr{\mbox{\PSyes{}}} میانگین \lr{\mbox{\ensuremath{\approx}}} واریانس در داده

\PSsubsection{زمان عدم استفاده}

\lr{\mbox{\PSno{}}} رویدادها مستقل نیستند (ورودهای خوشه‌ای)\PSbr{}\lr{\mbox{\PSno{}}} نرخ با زمان تغییر می‌کند (فرایند ناهمگن)\PSbr{}\lr{\mbox{\PSno{}}} رویدادها نادر نیستند (\lr{$p$} کوچک نیست) \lr{—} از دوجمله‌ای استفاده کنید\PSbr{}\lr{\mbox{\PSno{}}} میانگین \lr{\mbox{\ensuremath{\neq}}} واریانس (به بیش‌پراکندگی/کم‌پراکندگی توجه کنید)

\PSsubsection{پیاده‌سازی در \lr{\mbox{MATLAB}}}

\tcbset{PScodemode/.style={unbreakable}}

\begin{Shaded}
\begin{Highlighting}[]
\CommentTok{\%\% Poisson Distribution: Network Packet Arrivals}
\VariableTok{lambda} \OperatorTok{=} \FloatTok{4.5}\OperatorTok{;}     \CommentTok{\% Average packets per millisecond}
\VariableTok{k} \OperatorTok{=} \FloatTok{0}\OperatorTok{:}\FloatTok{15}\OperatorTok{;}

\VariableTok{pmf\_pois} \OperatorTok{=} \VariableTok{poisspdf}\NormalTok{(}\VariableTok{k}\OperatorTok{,} \VariableTok{lambda}\NormalTok{)}\OperatorTok{;}
\VariableTok{cdf\_pois} \OperatorTok{=} \VariableTok{poisscdf}\NormalTok{(}\VariableTok{k}\OperatorTok{,} \VariableTok{lambda}\NormalTok{)}\OperatorTok{;}

\VariableTok{figure}\OperatorTok{;}
\VariableTok{subplot}\NormalTok{(}\FloatTok{1}\OperatorTok{,}\FloatTok{2}\OperatorTok{,}\FloatTok{1}\NormalTok{)}\OperatorTok{;}
\VariableTok{stem}\NormalTok{(}\VariableTok{k}\OperatorTok{,} \VariableTok{pmf\_pois}\OperatorTok{,} \SpecialStringTok{\textquotesingle{}filled\textquotesingle{}}\OperatorTok{,} \SpecialStringTok{\textquotesingle{}LineWidth\textquotesingle{}}\OperatorTok{,} \FloatTok{1.5}\NormalTok{)}\OperatorTok{;}
\VariableTok{xlabel}\NormalTok{(}\SpecialStringTok{\textquotesingle{}Number of Arrivals (k)\textquotesingle{}}\NormalTok{)}\OperatorTok{;}
\VariableTok{ylabel}\NormalTok{(}\SpecialStringTok{\textquotesingle{}P(X = k)\textquotesingle{}}\NormalTok{)}\OperatorTok{;}
\VariableTok{title}\NormalTok{(}\VariableTok{sprintf}\NormalTok{(}\SpecialStringTok{\textquotesingle{}Poisson PMF (λ = \%.1f)\textquotesingle{}}\OperatorTok{,} \VariableTok{lambda}\NormalTok{))}\OperatorTok{;}
\VariableTok{grid} \VariableTok{on}\OperatorTok{;}

\VariableTok{subplot}\NormalTok{(}\FloatTok{1}\OperatorTok{,}\FloatTok{2}\OperatorTok{,}\FloatTok{2}\NormalTok{)}\OperatorTok{;}
\VariableTok{stairs}\NormalTok{(}\VariableTok{k}\OperatorTok{,} \VariableTok{cdf\_pois}\OperatorTok{,} \SpecialStringTok{\textquotesingle{}LineWidth\textquotesingle{}}\OperatorTok{,} \FloatTok{2}\NormalTok{)}\OperatorTok{;}
\VariableTok{xlabel}\NormalTok{(}\SpecialStringTok{\textquotesingle{}k\textquotesingle{}}\NormalTok{)}\OperatorTok{;} \VariableTok{ylabel}\NormalTok{(}\SpecialStringTok{\textquotesingle{}P(X ≤ k)\textquotesingle{}}\NormalTok{)}\OperatorTok{;}
\VariableTok{title}\NormalTok{(}\SpecialStringTok{\textquotesingle{}Poisson CDF\textquotesingle{}}\NormalTok{)}\OperatorTok{;}
\VariableTok{grid} \VariableTok{on}\OperatorTok{;}
\end{Highlighting}
\end{Shaded}

\begin{PSbox}[unbreakable]{ex}{مثال شبیه‌سازی: شمارش فوتون}

\tcbset{PScodemode/.style={unbreakable}}

\begin{Shaded}
\begin{Highlighting}[]
\CommentTok{\%\% Engineering Scenario: Photon Detection}
\CommentTok{\% An optical receiver detects photons. Average rate: 10 photons/μs.}
\CommentTok{\% What is P(detecting fewer than 5 photons in one μs)?}

\VariableTok{lambda} \OperatorTok{=} \FloatTok{10}\OperatorTok{;}   \CommentTok{\% average photons per μs}
\VariableTok{N} \OperatorTok{=} \FloatTok{100000}\OperatorTok{;}

\CommentTok{\% Analytical}
\VariableTok{p\_fewer\_5} \OperatorTok{=} \VariableTok{poisscdf}\NormalTok{(}\FloatTok{4}\OperatorTok{,} \VariableTok{lambda}\NormalTok{)}\OperatorTok{;}
\VariableTok{fprintf}\NormalTok{(}\SpecialStringTok{\textquotesingle{}P(X \textless{} 5): Analytical = \%.6f\textbackslash{}n\textquotesingle{}}\OperatorTok{,} \VariableTok{p\_fewer\_5}\NormalTok{)}\OperatorTok{;}

\CommentTok{\% Simulation}
\VariableTok{X\_sim} \OperatorTok{=} \VariableTok{poissrnd}\NormalTok{(}\VariableTok{lambda}\OperatorTok{,} \FloatTok{1}\OperatorTok{,} \VariableTok{N}\NormalTok{)}\OperatorTok{;}
\VariableTok{p\_fewer\_5\_sim} \OperatorTok{=} \VariableTok{mean}\NormalTok{(}\VariableTok{X\_sim} \OperatorTok{\textless{}} \FloatTok{5}\NormalTok{)}\OperatorTok{;}
\VariableTok{fprintf}\NormalTok{(}\SpecialStringTok{\textquotesingle{}P(X \textless{} 5): Simulation = \%.6f\textbackslash{}n\textquotesingle{}}\OperatorTok{,} \VariableTok{p\_fewer\_5\_sim}\NormalTok{)}\OperatorTok{;}

\CommentTok{\% Verify mean = variance (Poisson signature)}
\VariableTok{fprintf}\NormalTok{(}\SpecialStringTok{\textquotesingle{}Sample Mean = \%.3f, Sample Variance = \%.3f\textbackslash{}n\textquotesingle{}}\OperatorTok{,} \VariableTok{mean}\NormalTok{(}\VariableTok{X\_sim}\NormalTok{)}\OperatorTok{,} \VariableTok{var}\NormalTok{(}\VariableTok{X\_sim}\NormalTok{))}\OperatorTok{;}
\end{Highlighting}
\end{Shaded}

\end{PSbox}

\PSsection{3.7}{توزیع یکنواخت (پیوسته)}

\PSsubsection{تفسیر فیزیکی}

کمیتی را مدل می‌کند که \textbf{احتمال برابر} برای اختیار کردن هر مقدار در بازه \lr{$[a,\, b]$} دارد. نماینده \textbf{بیشترین عدم‌قطعیت} در یک محدوده معلوم است.

\PSsubsection{زمینه مهندسی}

\begin{itemize}
\tightlist
\item
  فاز یک سیگنال دریافتی (نامعلوم، با احتمال برابر در \lr{$[0,\, 2\pi)$})
\item
  خطای گِردکردن در یک مبدل آنالوگ به دیجیتال (یکنواخت در \lr{$\left[-\frac{\Delta}{2},\, \frac{\Delta}{2}\right]$} که \lr{$\Delta$} = گام کوانتیزاسیون)
\item
  زمان ورود در یک شکاف زمانی (موقعیت نامعلوم)
\item
  پروتکل دسترسی تصادفی: انتخاب یک زمان عقب‌نشینی تصادفی
\end{itemize}

\begin{PSbox}[unbreakable]{def}{تعریف ریاضی}

\lr{$X \sim \mathrm{Uniform}(a,\, b)$} یا \lr{$X \sim U(a,\, b)$}

\end{PSbox}

\PSsubsection{\lr{\mbox{PDF}}}

\PSdisplay{f_X(x) = \begin{cases} \frac{1}{b-a} & a \leq x \leq b \\ 0 & \PSmtext{otherwise} \end{cases}}

\PSsubsection{\lr{\mbox{CDF}}}

\PSdisplay{F_X(x) = \begin{cases} 0 & x < a \\ \frac{x-a}{b-a} & a \leq x \leq b \\ 1 & x > b \end{cases}}

\PSsubsection{میانگین}

\PSdisplay{E[X] = \frac{a+b}{2}}

\PSsubsection{واریانس}

\PSdisplay{\PSmtext{Var}(X) = \frac{(b-a)^2}{12}}

\PSsubsection{زمان استفاده}

\lr{\mbox{\PSyes{}}} تمام مقادیر در یک محدوده احتمال برابر دارند\PSbr{}\lr{\mbox{\PSyes{}}} هیچ اطلاعاتی به نفع مقدار خاصی نیست\PSbr{}\lr{\mbox{\PSyes{}}} مدل‌سازی خطای کوانتیزاسیون\PSbr{}\lr{\mbox{\PSyes{}}} فاز حامل بدون مدولاسیون

\PSsubsection{زمان عدم استفاده}

\lr{\mbox{\PSno{}}} مقادیر نزدیک به مرکز محتمل‌تر هستند (از گاوسی استفاده کنید)\PSbr{}\lr{\mbox{\PSno{}}} مقادیر نامنفی با کاهش نمایی هستند (از نمایی استفاده کنید)\PSbr{}\lr{\mbox{\PSno{}}} محدوده کران‌دار نیست

\PSsubsection{پیاده‌سازی در \lr{\mbox{MATLAB}}}

\tcbset{PScodemode/.style={unbreakable}}

\begin{Shaded}
\begin{Highlighting}[]
\CommentTok{\%\% Uniform Distribution: Quantization Error}
\VariableTok{a} \OperatorTok{=} \OperatorTok{{-}}\FloatTok{0.5}\OperatorTok{;}  \VariableTok{b} \OperatorTok{=} \FloatTok{0.5}\OperatorTok{;}  \CommentTok{\% Quantization step = 1 LSB}
\VariableTok{N} \OperatorTok{=} \FloatTok{100000}\OperatorTok{;}

\VariableTok{X} \OperatorTok{=} \VariableTok{a} \OperatorTok{+}\NormalTok{ (}\VariableTok{b}\OperatorTok{{-}}\VariableTok{a}\NormalTok{)}\OperatorTok{*}\VariableTok{rand}\NormalTok{(}\FloatTok{1}\OperatorTok{,} \VariableTok{N}\NormalTok{)}\OperatorTok{;}  \CommentTok{\% or: unifrnd(a, b, 1, N)}

\VariableTok{figure}\OperatorTok{;}
\VariableTok{histogram}\NormalTok{(}\VariableTok{X}\OperatorTok{,} \FloatTok{50}\OperatorTok{,} \SpecialStringTok{\textquotesingle{}Normalization\textquotesingle{}}\OperatorTok{,} \SpecialStringTok{\textquotesingle{}pdf\textquotesingle{}}\NormalTok{)}\OperatorTok{;}
\VariableTok{hold} \VariableTok{on}\OperatorTok{;}
\VariableTok{plot}\NormalTok{([}\VariableTok{a} \VariableTok{a} \VariableTok{b} \VariableTok{b}\NormalTok{]}\OperatorTok{,}\NormalTok{ [}\FloatTok{0} \FloatTok{1}\OperatorTok{/}\NormalTok{(}\VariableTok{b}\OperatorTok{{-}}\VariableTok{a}\NormalTok{) }\FloatTok{1}\OperatorTok{/}\NormalTok{(}\VariableTok{b}\OperatorTok{{-}}\VariableTok{a}\NormalTok{) }\FloatTok{0}\NormalTok{]}\OperatorTok{,} \SpecialStringTok{\textquotesingle{}r{-}\textquotesingle{}}\OperatorTok{,} \SpecialStringTok{\textquotesingle{}LineWidth\textquotesingle{}}\OperatorTok{,} \FloatTok{2}\NormalTok{)}\OperatorTok{;}
\VariableTok{xlabel}\NormalTok{(}\SpecialStringTok{\textquotesingle{}Quantization Error\textquotesingle{}}\NormalTok{)}\OperatorTok{;} \VariableTok{ylabel}\NormalTok{(}\SpecialStringTok{\textquotesingle{}f\_X(x)\textquotesingle{}}\NormalTok{)}\OperatorTok{;}
\VariableTok{title}\NormalTok{(}\SpecialStringTok{\textquotesingle{}Uniform PDF: Quantization Error\textquotesingle{}}\NormalTok{)}\OperatorTok{;}
\VariableTok{legend}\NormalTok{(}\SpecialStringTok{\textquotesingle{}Simulated\textquotesingle{}}\OperatorTok{,} \SpecialStringTok{\textquotesingle{}Theoretical\textquotesingle{}}\NormalTok{)}\OperatorTok{;}

\VariableTok{fprintf}\NormalTok{(}\SpecialStringTok{\textquotesingle{}Mean: Theory=\%.3f, Sim=\%.3f\textbackslash{}n\textquotesingle{}}\OperatorTok{,}\NormalTok{ (}\VariableTok{a}\OperatorTok{+}\VariableTok{b}\NormalTok{)}\OperatorTok{/}\FloatTok{2}\OperatorTok{,} \VariableTok{mean}\NormalTok{(}\VariableTok{X}\NormalTok{))}\OperatorTok{;}
\VariableTok{fprintf}\NormalTok{(}\SpecialStringTok{\textquotesingle{}Var:  Theory=\%.4f, Sim=\%.4f\textbackslash{}n\textquotesingle{}}\OperatorTok{,}\NormalTok{ (}\VariableTok{b}\OperatorTok{{-}}\VariableTok{a}\NormalTok{)}\OperatorTok{\^{}}\FloatTok{2}\OperatorTok{/}\FloatTok{12}\OperatorTok{,} \VariableTok{var}\NormalTok{(}\VariableTok{X}\NormalTok{))}\OperatorTok{;}
\end{Highlighting}
\end{Shaded}

\PSsection{3.8}{توزیع نمایی}

\PSsubsection{تفسیر فیزیکی}

\textbf{زمان میان رویدادها} در یک فرایند پواسون، یا \textbf{عمر} یک قطعه بی‌حافظه را مدل می‌کند.

\PSsubsection{زمینه مهندسی}

\begin{itemize}
\tightlist
\item
  زمان میان ورود بسته‌ها
\item
  زمان تا خرابی قطعه (نرخ خرابی ثابت)
\item
  زمان میان رویدادهای پرتو کیهانی در یک آشکارساز
\item
  زمان بین‌ورودی تماس‌ها در یک مرکز سوییچ
\item
  مدت یک فروافت محوشدگی در یک کانال بی‌سیم
\end{itemize}

\begin{PSbox}[unbreakable]{def}{تعریف ریاضی}

\lr{$X \sim \mathrm{Exponential}(\lambda)$} که در آن \lr{$\lambda$} = پارامتر نرخ (رویداد در واحد زمان).

پارامتربندی جایگزین: \lr{$X \sim \mathrm{Exp}(\beta)$} که در آن \lr{$\beta = \frac{1}{\lambda}$} = میانگین.

\end{PSbox}

\PSsubsection{\lr{\mbox{PDF}}}

\PSdisplay{f_X(x) = \lambda e^{-\lambda x}, \quad x \geq 0}

یا با میانگین \lr{\mbox{$\beta$: $f_{X}(x) = \frac{1}{\beta}e^{-x/\beta}$}}

\PSsubsection{\lr{\mbox{CDF}}}

\PSdisplay{F_X(x) = 1 - e^{-\lambda x}, \quad x \geq 0}

\PSsubsection{میانگین}

\PSdisplay{E[X] = \frac{1}{\lambda} = \beta}

\PSsubsection{واریانس}

\PSdisplay{\PSmtext{Var}(X) = \frac{1}{\lambda^2} = \beta^2}

\textbf{توجه:} برای توزیع نمایی، انحراف معیار برابر میانگین است.

\PSsubsection{خاصیت بی‌حافظه}

\PSdisplay{P(X > s + t \mid X > s) = P(X > t)}

توزیع نمایی تنها توزیع پیوسته با این خاصیت است.

\textbf{معنای مهندسی:} قطعه‌ای که \lr{$s$} ساعت کار کرده است، از نظر آماری «به‌خوبی نو» است. عمر باقی‌مانده آن همان توزیع یک قطعه کاملاً نو را دارد.

\PSsubsection{رابطه با پواسون}

اگر رویدادها به‌صورت یک فرایند پواسون با نرخ \lr{$\lambda$} وارد شوند:

\begin{itemize}
\tightlist
\item
  تعداد رویدادها در زمان \lr{$T \sim \mathrm{Poisson}(\lambda T)$}
\item
  زمان میان رویدادها \lr{$\sim \mathrm{Exponential}(\lambda)$}
\end{itemize}

\PSsubsection{زمان استفاده}

\lr{\mbox{\PSyes{}}} مدل‌سازی زمان میان رویدادها در یک فرایند پواسون\PSbr{}\lr{\mbox{\PSyes{}}} عمر قطعه با نرخ خرابی ثابت\PSbr{}\lr{\mbox{\PSyes{}}} فرض بی‌حافظه معقول است

\PSsubsection{زمان عدم استفاده}

\lr{\mbox{\PSno{}}} نرخ خرابی با سن افزایش می‌یابد (فرسودگی) \lr{—} از وایبول یا گاما استفاده کنید\PSbr{}\lr{\mbox{\PSno{}}} نرخ خرابی با سن کاهش می‌یابد (سوختگی اولیه) \lr{—} از وایبول استفاده کنید\PSbr{}\lr{\mbox{\PSno{}}} داده افزایش خطر را نشان می‌دهد \lr{—} بی‌حافظه نیست

\PSsubsection{پیاده‌سازی در \lr{\mbox{MATLAB}}}

\tcbset{PScodemode/.style={unbreakable}}

\begin{Shaded}
\begin{Highlighting}[]
\CommentTok{\%\% Exponential Distribution: Component Lifetime}
\VariableTok{lambda} \OperatorTok{=} \FloatTok{0.001}\OperatorTok{;}          \CommentTok{\% Failure rate: 0.001 per hour}
\VariableTok{beta} \OperatorTok{=} \FloatTok{1}\OperatorTok{/}\VariableTok{lambda}\OperatorTok{;}         \CommentTok{\% Mean lifetime: 1000 hours}

\VariableTok{x} \OperatorTok{=} \VariableTok{linspace}\NormalTok{(}\FloatTok{0}\OperatorTok{,} \FloatTok{5000}\OperatorTok{,} \FloatTok{1000}\NormalTok{)}\OperatorTok{;}
\VariableTok{pdf\_exp} \OperatorTok{=} \VariableTok{exppdf}\NormalTok{(}\VariableTok{x}\OperatorTok{,} \VariableTok{beta}\NormalTok{)}\OperatorTok{;}
\VariableTok{cdf\_exp} \OperatorTok{=} \VariableTok{expcdf}\NormalTok{(}\VariableTok{x}\OperatorTok{,} \VariableTok{beta}\NormalTok{)}\OperatorTok{;}

\VariableTok{figure}\OperatorTok{;}
\VariableTok{subplot}\NormalTok{(}\FloatTok{2}\OperatorTok{,}\FloatTok{1}\OperatorTok{,}\FloatTok{1}\NormalTok{)}\OperatorTok{;}
\VariableTok{plot}\NormalTok{(}\VariableTok{x}\OperatorTok{,} \VariableTok{pdf\_exp}\OperatorTok{,} \SpecialStringTok{\textquotesingle{}b{-}\textquotesingle{}}\OperatorTok{,} \SpecialStringTok{\textquotesingle{}LineWidth\textquotesingle{}}\OperatorTok{,} \FloatTok{2}\NormalTok{)}\OperatorTok{;}
\VariableTok{xlabel}\NormalTok{(}\SpecialStringTok{\textquotesingle{}Time (hours)\textquotesingle{}}\NormalTok{)}\OperatorTok{;} \VariableTok{ylabel}\NormalTok{(}\SpecialStringTok{\textquotesingle{}f\_X(x)\textquotesingle{}}\NormalTok{)}\OperatorTok{;}
\VariableTok{title}\NormalTok{(}\VariableTok{sprintf}\NormalTok{(}\SpecialStringTok{\textquotesingle{}Exponential PDF: Mean Lifetime = \%d hours\textquotesingle{}}\OperatorTok{,} \VariableTok{beta}\NormalTok{))}\OperatorTok{;}
\VariableTok{grid} \VariableTok{on}\OperatorTok{;}

\VariableTok{subplot}\NormalTok{(}\FloatTok{2}\OperatorTok{,}\FloatTok{1}\OperatorTok{,}\FloatTok{2}\NormalTok{)}\OperatorTok{;}
\VariableTok{plot}\NormalTok{(}\VariableTok{x}\OperatorTok{,} \VariableTok{cdf\_exp}\OperatorTok{,} \SpecialStringTok{\textquotesingle{}r{-}\textquotesingle{}}\OperatorTok{,} \SpecialStringTok{\textquotesingle{}LineWidth\textquotesingle{}}\OperatorTok{,} \FloatTok{2}\NormalTok{)}\OperatorTok{;}
\VariableTok{xlabel}\NormalTok{(}\SpecialStringTok{\textquotesingle{}Time (hours)\textquotesingle{}}\NormalTok{)}\OperatorTok{;} \VariableTok{ylabel}\NormalTok{(}\SpecialStringTok{\textquotesingle{}F\_X(x) = P(failure by time x)\textquotesingle{}}\NormalTok{)}\OperatorTok{;}
\VariableTok{title}\NormalTok{(}\SpecialStringTok{\textquotesingle{}Exponential CDF\textquotesingle{}}\NormalTok{)}\OperatorTok{;}
\VariableTok{grid} \VariableTok{on}\OperatorTok{;}

\CommentTok{\% P(component lasts more than 2000 hours)}
\VariableTok{p\_survive\_2000} \OperatorTok{=} \FloatTok{1} \OperatorTok{{-}} \VariableTok{expcdf}\NormalTok{(}\FloatTok{2000}\OperatorTok{,} \VariableTok{beta}\NormalTok{)}\OperatorTok{;}
\VariableTok{fprintf}\NormalTok{(}\SpecialStringTok{\textquotesingle{}P(T \textgreater{} 2000 hours) = \%.4f\textbackslash{}n\textquotesingle{}}\OperatorTok{,} \VariableTok{p\_survive\_2000}\NormalTok{)}\OperatorTok{;}
\end{Highlighting}
\end{Shaded}

\begin{PSbox}[unbreakable]{ex}{مثال شبیه‌سازی: بررسی خاصیت بی‌حافظه}

\tcbset{PScodemode/.style={unbreakable}}

\begin{Shaded}
\begin{Highlighting}[]
\CommentTok{\%\% Verify Memoryless Property}
\VariableTok{beta} \OperatorTok{=} \FloatTok{100}\OperatorTok{;}   \CommentTok{\% Mean lifetime = 100 hours}
\VariableTok{N} \OperatorTok{=} \FloatTok{100000}\OperatorTok{;}

\VariableTok{X} \OperatorTok{=} \VariableTok{exprnd}\NormalTok{(}\VariableTok{beta}\OperatorTok{,} \FloatTok{1}\OperatorTok{,} \VariableTok{N}\NormalTok{)}\OperatorTok{;}

\CommentTok{\% Conditional: given X \textgreater{} 50, what is distribution of (X {-} 50)?}
\VariableTok{survivors} \OperatorTok{=} \VariableTok{X}\NormalTok{(}\VariableTok{X} \OperatorTok{\textgreater{}} \FloatTok{50}\NormalTok{)}\OperatorTok{;}
\VariableTok{remaining\_life} \OperatorTok{=} \VariableTok{survivors} \OperatorTok{{-}} \FloatTok{50}\OperatorTok{;}

\VariableTok{fprintf}\NormalTok{(}\SpecialStringTok{\textquotesingle{}Mean of X: \%.1f (should be \%.1f)\textbackslash{}n\textquotesingle{}}\OperatorTok{,} \VariableTok{mean}\NormalTok{(}\VariableTok{X}\NormalTok{)}\OperatorTok{,} \VariableTok{beta}\NormalTok{)}\OperatorTok{;}
\VariableTok{fprintf}\NormalTok{(}\SpecialStringTok{\textquotesingle{}Mean of remaining life given survival past 50: \%.1f\textbackslash{}n\textquotesingle{}}\OperatorTok{,} \VariableTok{mean}\NormalTok{(}\VariableTok{remaining\_life}\NormalTok{))}\OperatorTok{;}
\VariableTok{fprintf}\NormalTok{(}\SpecialStringTok{\textquotesingle{}Should also be \%.1f (memoryless!)\textbackslash{}n\textquotesingle{}}\OperatorTok{,} \VariableTok{beta}\NormalTok{)}\OperatorTok{;}
\end{Highlighting}
\end{Shaded}

\end{PSbox}

\PSsection{3.9}{توزیع گاوسی (نرمال)}

\PSsubsection{تفسیر فیزیکی}

کمیت‌هایی را مدل می‌کند که از \textbf{جمع بسیاری از اثرهای تصادفی کوچک و مستقل} پدید می‌آیند. به‌دلیل قضیه حد مرکزی، مهم‌ترین توزیع در مهندسی است.

\PSsubsection{زمینه مهندسی}

\begin{itemize}
\tightlist
\item
  ولتاژ نویز حرارتی (جمع سهم‌های بسیاری از الکترون‌ها)
\item
  خطاهای اندازه‌گیری (جمع بسیاری از منابع خطای کوچک)
\item
  تغییرات ساخت (بسیاری از تغییرات مستقل فرایند)
\item
  تداخل تجمعی در مخابرات
\item
  هر کمیتی که به‌خوبی با \lr{\mbox{CLT}} مدل می‌شود
\end{itemize}

\begin{PSbox}[unbreakable]{def}{تعریف ریاضی}

\lr{$X \sim N(\mu,\, \sigma^{2})$} که در آن \lr{$\mu$} = میانگین و \lr{$\sigma^{2}$} = واریانس.

\end{PSbox}

\PSsubsection{\lr{\mbox{PDF}}}

\PSdisplay{f_X(x) = \frac{1}{\sigma\sqrt{2\pi}} \exp\left(-\frac{(x-\mu)^2}{2\sigma^2}\right), \quad -\infty < x < \infty}

\PSsubsection{\lr{\mbox{CDF}}}

\PSdisplay{F_X(x) = \Phi\left(\frac{x-\mu}{\sigma}\right) = \frac{1}{2}\left[1 + \PSmtext{erf}\left(\frac{x-\mu}{\sigma\sqrt{2}}\right)\right]}

فرم بسته ندارد \lr{—} به‌صورت عددی یا با جدول‌ها محاسبه می‌شود.

\PSsubsection{میانگین}

\PSdisplay{E[X] = \mu}

\PSsubsection{واریانس}

\PSdisplay{\PSmtext{Var}(X) = \sigma^2}

\PSsubsection{نرمال استاندارد}

\lr{$Z \sim N(0,\, 1)$}: شکل استانداردشده.

هر توزیع نرمال را می‌توان استانداردسازی کرد: \lr{$Z = \frac{X - \mu}{\sigma}$}

\PSsubsection{قواعد کلیدی احتمال (قواعد \lr{$\sigma$})}

{\def\LTcaptype{none} % do not increment counter
\begin{longtable}[c]{@{}cc@{}}
\toprule\noalign{}
\rowcolor{PSnavy}\PSth{محدوده} & \PSth{احتمال} \\
\midrule\noalign{}
\endhead
\bottomrule\noalign{}
\endlastfoot
\lr{$\mu \pm 1\sigma$} & \PSnum{68.27\%} \\
\lr{$\mu \pm 2\sigma$} & \PSnum{95.45\%} \\
\lr{$\mu \pm 3\sigma$} & \PSnum{99.73\%} \\
\lr{$\mu \pm 4\sigma$} & \PSnum{99.9937\%} \\
\end{longtable}
}

\PSsubsection{ویژگی‌ها}

\begin{enumerate}
\def\labelenumi{\arabic{enumi}.}
\tightlist
\item
  \textbf{متقارن} حول \lr{$\mu$}
\item
  \textbf{تبدیل خطی:} اگر \lr{$X \sim N(\mu,\, \sigma^{2})$}، آنگاه \lr{\mbox{aX $+ b \sim N(a\mu + b,\, a^{2}\sigma^{2})$}}
\item
  \textbf{جمع گاوسی‌ها:} اگر \lr{$X \sim N(\mu_{1},\, \sigma_{1}^{2})$} و \lr{$Y \sim N(\mu_{2},\, \sigma_{2}^{2})$} مستقل باشند، آنگاه \lr{$X + Y \sim N(\mu_{1}+\mu_{2},\, \sigma_{1}^{2}+\sigma_{2}^{2})$}
\item
  \textbf{کاملاً مشخص‌شده} با میانگین و واریانس (تمام گشتاورهای مراتب بالاتر تعیین می‌شوند)
\end{enumerate}

\PSsubsection{زمان استفاده}

\lr{\mbox{\PSyes{}}} نویز حرارتی در مدارها\PSbr{}\lr{\mbox{\PSyes{}}} خطاهای اندازه‌گیری\PSbr{}\lr{\mbox{\PSyes{}}} هر جمعی از بسیاری از اثرهای تصادفی مستقل \lr{\mbox{(CLT)}}\PSbr{}\lr{\mbox{\PSyes{}}} تغییرات ساخت\PSbr{}\lr{\mbox{\PSyes{}}} هنگامی که هیستوگرام داده زنگوله‌ای و متقارن است

\PSsubsection{زمان عدم استفاده}

\lr{\mbox{\PSno{}}} کمیت‌های اکیداً مثبت (گاوسی مقادیر منفی را مجاز می‌داند)\PSbr{}\lr{\mbox{\PSno{}}} داده‌های به‌شدت چوله\PSbr{}\lr{\mbox{\PSno{}}} کمیت‌های کران‌دار (گاوسی تکیه‌گاه نامتناهی دارد)\PSbr{}\lr{\mbox{\PSno{}}} شمارش‌های گسسته\PSbr{}\lr{\mbox{\PSno{}}} پدیده‌های سنگین‌دنباله (از توزیع \lr{$t$} یا کوشی استفاده کنید)

\PSsubsection{پیاده‌سازی در \lr{\mbox{MATLAB}}}

\tcbset{PScodemode/.style={unbreakable}}

\begin{Shaded}
\begin{Highlighting}[]
\CommentTok{\%\% Gaussian Distribution: Thermal Noise}
\VariableTok{mu} \OperatorTok{=} \FloatTok{0}\OperatorTok{;}           \CommentTok{\% Zero{-}mean noise}
\VariableTok{sigma} \OperatorTok{=} \FloatTok{0.1}\OperatorTok{;}      \CommentTok{\% RMS noise = 100 mV}

\VariableTok{x} \OperatorTok{=} \VariableTok{linspace}\NormalTok{(}\VariableTok{mu} \OperatorTok{{-}} \FloatTok{4}\OperatorTok{*}\VariableTok{sigma}\OperatorTok{,} \VariableTok{mu} \OperatorTok{+} \FloatTok{4}\OperatorTok{*}\VariableTok{sigma}\OperatorTok{,} \FloatTok{1000}\NormalTok{)}\OperatorTok{;}
\VariableTok{pdf\_norm} \OperatorTok{=} \VariableTok{normpdf}\NormalTok{(}\VariableTok{x}\OperatorTok{,} \VariableTok{mu}\OperatorTok{,} \VariableTok{sigma}\NormalTok{)}\OperatorTok{;}

\VariableTok{figure}\OperatorTok{;}
\VariableTok{plot}\NormalTok{(}\VariableTok{x}\OperatorTok{,} \VariableTok{pdf\_norm}\OperatorTok{,} \SpecialStringTok{\textquotesingle{}b{-}\textquotesingle{}}\OperatorTok{,} \SpecialStringTok{\textquotesingle{}LineWidth\textquotesingle{}}\OperatorTok{,} \FloatTok{2}\NormalTok{)}\OperatorTok{;}
\VariableTok{hold} \VariableTok{on}\OperatorTok{;}
\CommentTok{\% Shade the ±2σ region}
\VariableTok{x\_fill} \OperatorTok{=} \VariableTok{linspace}\NormalTok{(}\VariableTok{mu}\OperatorTok{{-}}\FloatTok{2}\OperatorTok{*}\VariableTok{sigma}\OperatorTok{,} \VariableTok{mu}\OperatorTok{+}\FloatTok{2}\OperatorTok{*}\VariableTok{sigma}\OperatorTok{,} \FloatTok{200}\NormalTok{)}\OperatorTok{;}
\VariableTok{fill}\NormalTok{([}\VariableTok{x\_fill}\OperatorTok{,} \VariableTok{fliplr}\NormalTok{(}\VariableTok{x\_fill}\NormalTok{)]}\OperatorTok{,}\NormalTok{ [}\VariableTok{normpdf}\NormalTok{(}\VariableTok{x\_fill}\OperatorTok{,}\VariableTok{mu}\OperatorTok{,}\VariableTok{sigma}\NormalTok{)}\OperatorTok{,} \VariableTok{zeros}\NormalTok{(}\VariableTok{size}\NormalTok{(}\VariableTok{x\_fill}\NormalTok{))]}\OperatorTok{,} \OperatorTok{...}
     \SpecialStringTok{\textquotesingle{}b\textquotesingle{}}\OperatorTok{,} \SpecialStringTok{\textquotesingle{}FaceAlpha\textquotesingle{}}\OperatorTok{,} \FloatTok{0.3}\NormalTok{)}\OperatorTok{;}
\VariableTok{xlabel}\NormalTok{(}\SpecialStringTok{\textquotesingle{}Noise Voltage (V)\textquotesingle{}}\NormalTok{)}\OperatorTok{;}
\VariableTok{ylabel}\NormalTok{(}\SpecialStringTok{\textquotesingle{}f\_X(x)\textquotesingle{}}\NormalTok{)}\OperatorTok{;}
\VariableTok{title}\NormalTok{(}\VariableTok{sprintf}\NormalTok{(}\SpecialStringTok{\textquotesingle{}Gaussian PDF: μ=\%.1f, σ=\%.3f V\textquotesingle{}}\OperatorTok{,} \VariableTok{mu}\OperatorTok{,} \VariableTok{sigma}\NormalTok{))}\OperatorTok{;}
\VariableTok{legend}\NormalTok{(}\SpecialStringTok{\textquotesingle{}PDF\textquotesingle{}}\OperatorTok{,} \SpecialStringTok{\textquotesingle{}±2σ region (95.45\%)\textquotesingle{}}\NormalTok{)}\OperatorTok{;}
\VariableTok{grid} \VariableTok{on}\OperatorTok{;}

\CommentTok{\% Probability computations}
\VariableTok{fprintf}\NormalTok{(}\SpecialStringTok{\textquotesingle{}P(|noise| \textgreater{} 0.2V) = \%.4f\textbackslash{}n\textquotesingle{}}\OperatorTok{,} \FloatTok{2}\OperatorTok{*}\NormalTok{(}\FloatTok{1}\OperatorTok{{-}}\VariableTok{normcdf}\NormalTok{(}\FloatTok{0.2}\OperatorTok{,} \FloatTok{0}\OperatorTok{,} \VariableTok{sigma}\NormalTok{)))}\OperatorTok{;}
\VariableTok{fprintf}\NormalTok{(}\SpecialStringTok{\textquotesingle{}P(|noise| \textgreater{} 0.3V) = \%.4f\textbackslash{}n\textquotesingle{}}\OperatorTok{,} \FloatTok{2}\OperatorTok{*}\NormalTok{(}\FloatTok{1}\OperatorTok{{-}}\VariableTok{normcdf}\NormalTok{(}\FloatTok{0.3}\OperatorTok{,} \FloatTok{0}\OperatorTok{,} \VariableTok{sigma}\NormalTok{)))}\OperatorTok{;}
\end{Highlighting}
\end{Shaded}

\begin{PSbox}[breakable]{ex}{مثال شبیه‌سازی: تابع \lr{$Q$} و \lr{\mbox{BER}}}

\tcbset{PScodemode/.style={unbreakable}}

\begin{Shaded}
\begin{Highlighting}[]
\CommentTok{\%\% Engineering Application: Binary Communication BER}
\CommentTok{\% In BPSK with AWGN: BER = Q(sqrt(2*Eb/N0))}
\CommentTok{\% Q(x) = 1 {-} Phi(x) = probability that N(0,1) \textgreater{} x}

\VariableTok{EbN0\_dB} \OperatorTok{=} \FloatTok{0}\OperatorTok{:}\FloatTok{0.5}\OperatorTok{:}\FloatTok{12}\OperatorTok{;}
\VariableTok{EbN0} \OperatorTok{=} \FloatTok{10}\OperatorTok{.\^{}}\NormalTok{(}\VariableTok{EbN0\_dB}\OperatorTok{/}\FloatTok{10}\NormalTok{)}\OperatorTok{;}

\CommentTok{\% Theoretical BER}
\VariableTok{BER\_theory} \OperatorTok{=} \VariableTok{qfunc}\NormalTok{(}\VariableTok{sqrt}\NormalTok{(}\FloatTok{2}\OperatorTok{*}\VariableTok{EbN0}\NormalTok{))}\OperatorTok{;}

\CommentTok{\% Simulated BER}
\VariableTok{N\_bits} \OperatorTok{=} \FloatTok{1e6}\OperatorTok{;}
\VariableTok{BER\_sim} \OperatorTok{=} \VariableTok{zeros}\NormalTok{(}\VariableTok{size}\NormalTok{(}\VariableTok{EbN0\_dB}\NormalTok{))}\OperatorTok{;}
\KeywordTok{for} \VariableTok{idx} \OperatorTok{=} \FloatTok{1}\OperatorTok{:}\VariableTok{length}\NormalTok{(}\VariableTok{EbN0}\NormalTok{)}
    \CommentTok{\% BPSK: transmit ±1}
    \VariableTok{bits} \OperatorTok{=} \FloatTok{2}\OperatorTok{*}\VariableTok{randi}\NormalTok{([}\FloatTok{0}\OperatorTok{,}\FloatTok{1}\NormalTok{]}\OperatorTok{,} \FloatTok{1}\OperatorTok{,} \VariableTok{N\_bits}\NormalTok{) }\OperatorTok{{-}} \FloatTok{1}\OperatorTok{;}
    \VariableTok{noise\_std} \OperatorTok{=} \FloatTok{1}\OperatorTok{/}\VariableTok{sqrt}\NormalTok{(}\FloatTok{2}\OperatorTok{*}\VariableTok{EbN0}\NormalTok{(}\VariableTok{idx}\NormalTok{))}\OperatorTok{;}
    \VariableTok{received} \OperatorTok{=} \VariableTok{bits} \OperatorTok{+} \VariableTok{noise\_std} \OperatorTok{*} \VariableTok{randn}\NormalTok{(}\FloatTok{1}\OperatorTok{,} \VariableTok{N\_bits}\NormalTok{)}\OperatorTok{;}
    \VariableTok{decisions} \OperatorTok{=} \VariableTok{sign}\NormalTok{(}\VariableTok{received}\NormalTok{)}\OperatorTok{;}
    \VariableTok{BER\_sim}\NormalTok{(}\VariableTok{idx}\NormalTok{) }\OperatorTok{=} \VariableTok{mean}\NormalTok{(}\VariableTok{decisions} \OperatorTok{\textasciitilde{}=} \VariableTok{bits}\NormalTok{)}\OperatorTok{;}
\KeywordTok{end}

\VariableTok{figure}\OperatorTok{;}
\VariableTok{semilogy}\NormalTok{(}\VariableTok{EbN0\_dB}\OperatorTok{,} \VariableTok{BER\_theory}\OperatorTok{,} \SpecialStringTok{\textquotesingle{}b{-}\textquotesingle{}}\OperatorTok{,} \SpecialStringTok{\textquotesingle{}LineWidth\textquotesingle{}}\OperatorTok{,} \FloatTok{2}\NormalTok{)}\OperatorTok{;}
\VariableTok{hold} \VariableTok{on}\OperatorTok{;}
\VariableTok{semilogy}\NormalTok{(}\VariableTok{EbN0\_dB}\OperatorTok{,} \VariableTok{BER\_sim}\OperatorTok{,} \SpecialStringTok{\textquotesingle{}ro\textquotesingle{}}\OperatorTok{,} \SpecialStringTok{\textquotesingle{}MarkerSize\textquotesingle{}}\OperatorTok{,} \FloatTok{6}\NormalTok{)}\OperatorTok{;}
\VariableTok{xlabel}\NormalTok{(}\SpecialStringTok{\textquotesingle{}E\_b/N\_0 (dB)\textquotesingle{}}\NormalTok{)}\OperatorTok{;} \VariableTok{ylabel}\NormalTok{(}\SpecialStringTok{\textquotesingle{}Bit Error Rate\textquotesingle{}}\NormalTok{)}\OperatorTok{;}
\VariableTok{title}\NormalTok{(}\SpecialStringTok{\textquotesingle{}BPSK BER: Theory vs Simulation\textquotesingle{}}\NormalTok{)}\OperatorTok{;}
\VariableTok{legend}\NormalTok{(}\SpecialStringTok{\textquotesingle{}Theory: Q(√(2E\_b/N\_0))\textquotesingle{}}\OperatorTok{,} \SpecialStringTok{\textquotesingle{}Simulation\textquotesingle{}}\NormalTok{)}\OperatorTok{;}
\VariableTok{grid} \VariableTok{on}\OperatorTok{;}
\end{Highlighting}
\end{Shaded}

\end{PSbox}

\PSsection{3.10}{توزیع گاما}

\PSsubsection{تفسیر فیزیکی}

\textbf{کل زمان انتظار تا \lr{\mbox{k}}-امین رویداد} در یک فرایند پواسون، یا هر جمعی از \lr{$k$} متغیر تصادفی نمایی مستقل را مدل می‌کند.

\PSsubsection{زمینه مهندسی}

\begin{itemize}
\tightlist
\item
  کل زمان تعمیر برای \lr{$k$} قطعه (هر یک با زمان تعمیر نمایی)
\item
  زمان تا \lr{\mbox{k}}-امین ورود بسته
\item
  زمان سرویس تجمعی برای \lr{$k$} کار صف‌بندی‌شده
\item
  مدل‌سازی مقدار بارش
\end{itemize}

\begin{PSbox}[unbreakable]{def}{تعریف ریاضی}

\lr{$X \sim \mathrm{Gamma}(\alpha,\, \beta)$} که در آن \lr{$\alpha$} = پارامتر شکل و \lr{$\beta$} = پارامتر مقیاس (یا نرخ \lr{$\lambda = \frac{1}{\beta}$}).

\end{PSbox}

\PSsubsection{\lr{\mbox{PDF}}}

\PSdisplay{f_X(x) = \frac{x^{\alpha-1} e^{-x/\beta}}{\beta^\alpha \Gamma(\alpha)}, \quad x \geq 0}

که در آن برای \lr{$\alpha$} صحیح \lr{$\Gamma (\alpha) = (\alpha -1)$}!، و در حالت کلی \lr{$\Gamma (\alpha) = \int_{0}^{\infty} t^{\alpha -1}e^{-t}dt$}.

\PSsubsection{\lr{\mbox{CDF}}}

برای \lr{$\alpha$} کلی فرم بسته ندارد. از طریق تابع گامای ناقص در دسترس است.

\PSsubsection{میانگین}

\PSdisplay{E[X] = \alpha\beta}

\PSsubsection{واریانس}

\PSdisplay{\PSmtext{Var}(X) = \alpha\beta^2}

\PSsubsection{حالت‌های خاص}

{\def\LTcaptype{none} % do not increment counter
\begin{longtable}[c]{@{}cc@{}}
\toprule\noalign{}
\rowcolor{PSnavy}\PSth{پارامترها} & \PSth{توزیع} \\
\midrule\noalign{}
\endhead
\bottomrule\noalign{}
\endlastfoot
\lr{$\alpha = 1$} & \lr{$\mathrm{Exponential}(\beta)$} \\
\lr{$\alpha = \frac{n}{2},\, \beta = 2$} & \lr{\mbox{Chi-square(n)}} \\
\lr{\mbox{$\alpha$ = integer $k$}} & \lr{$\mathrm{Erlang}(k,\, \beta)$} \\
\end{longtable}
}

\PSsubsection{زمان استفاده}

\lr{\mbox{\PSyes{}}} جمع عمرهای نمایی مستقل\PSbr{}\lr{\mbox{\PSyes{}}} زمان انتظار برای \lr{\mbox{k}}-امین رویداد پواسون\PSbr{}\lr{\mbox{\PSyes{}}} داده‌های چوله و نامنفی\PSbr{}\lr{\mbox{\PSyes{}}} شکل انعطاف‌پذیر (چولگی قابل تنظیم)

\PSsubsection{زمان عدم استفاده}

\lr{\mbox{\PSno{}}} داده می‌تواند منفی باشد\PSbr{}\lr{\mbox{\PSno{}}} داده متقارن است (از گاوسی استفاده کنید)\PSbr{}\lr{\mbox{\PSno{}}} توزیع ساده‌تری برازش می‌شود (نمایی برای \lr{$k = 1$})

\PSsubsection{پیاده‌سازی در \lr{\mbox{MATLAB}}}

\tcbset{PScodemode/.style={unbreakable}}

\begin{Shaded}
\begin{Highlighting}[]
\CommentTok{\%\% Gamma Distribution: Total Repair Time}
\CommentTok{\% 3 components must be repaired sequentially}
\CommentTok{\% Each repair time \textasciitilde{} Exp(mean = 2 hours)}
\VariableTok{alpha} \OperatorTok{=} \FloatTok{3}\OperatorTok{;}         \CommentTok{\% shape (number of stages)}
\VariableTok{beta} \OperatorTok{=} \FloatTok{2}\OperatorTok{;}          \CommentTok{\% scale (mean per stage)}

\VariableTok{x} \OperatorTok{=} \VariableTok{linspace}\NormalTok{(}\FloatTok{0}\OperatorTok{,} \FloatTok{20}\OperatorTok{,} \FloatTok{500}\NormalTok{)}\OperatorTok{;}
\VariableTok{pdf\_gamma} \OperatorTok{=} \VariableTok{gampdf}\NormalTok{(}\VariableTok{x}\OperatorTok{,} \VariableTok{alpha}\OperatorTok{,} \VariableTok{beta}\NormalTok{)}\OperatorTok{;}

\VariableTok{figure}\OperatorTok{;}
\VariableTok{plot}\NormalTok{(}\VariableTok{x}\OperatorTok{,} \VariableTok{pdf\_gamma}\OperatorTok{,} \SpecialStringTok{\textquotesingle{}b{-}\textquotesingle{}}\OperatorTok{,} \SpecialStringTok{\textquotesingle{}LineWidth\textquotesingle{}}\OperatorTok{,} \FloatTok{2}\NormalTok{)}\OperatorTok{;}
\VariableTok{xlabel}\NormalTok{(}\SpecialStringTok{\textquotesingle{}Total Repair Time (hours)\textquotesingle{}}\NormalTok{)}\OperatorTok{;}
\VariableTok{ylabel}\NormalTok{(}\SpecialStringTok{\textquotesingle{}f\_X(x)\textquotesingle{}}\NormalTok{)}\OperatorTok{;}
\VariableTok{title}\NormalTok{(}\VariableTok{sprintf}\NormalTok{(}\SpecialStringTok{\textquotesingle{}Gamma PDF: α=\%d, β=\%d\textquotesingle{}}\OperatorTok{,} \VariableTok{alpha}\OperatorTok{,} \VariableTok{beta}\NormalTok{))}\OperatorTok{;}
\VariableTok{grid} \VariableTok{on}\OperatorTok{;}

\CommentTok{\% Verify by summing exponentials}
\VariableTok{N} \OperatorTok{=} \FloatTok{100000}\OperatorTok{;}
\VariableTok{X\_sim} \OperatorTok{=} \VariableTok{sum}\NormalTok{(}\VariableTok{exprnd}\NormalTok{(}\VariableTok{beta}\OperatorTok{,} \VariableTok{alpha}\OperatorTok{,} \VariableTok{N}\NormalTok{)}\OperatorTok{,} \FloatTok{1}\NormalTok{)}\OperatorTok{;}  \CommentTok{\% Sum of 3 exponentials}
\VariableTok{fprintf}\NormalTok{(}\SpecialStringTok{\textquotesingle{}Mean: Theory=\%.2f, Sim=\%.2f\textbackslash{}n\textquotesingle{}}\OperatorTok{,} \VariableTok{alpha}\OperatorTok{*}\VariableTok{beta}\OperatorTok{,} \VariableTok{mean}\NormalTok{(}\VariableTok{X\_sim}\NormalTok{))}\OperatorTok{;}
\VariableTok{fprintf}\NormalTok{(}\SpecialStringTok{\textquotesingle{}Var:  Theory=\%.2f, Sim=\%.2f\textbackslash{}n\textquotesingle{}}\OperatorTok{,} \VariableTok{alpha}\OperatorTok{*}\VariableTok{beta}\OperatorTok{\^{}}\FloatTok{2}\OperatorTok{,} \VariableTok{var}\NormalTok{(}\VariableTok{X\_sim}\NormalTok{))}\OperatorTok{;}
\end{Highlighting}
\end{Shaded}

\PSsection{3.11}{توزیع رایلی}

\PSsubsection{تفسیر فیزیکی}

\textbf{اندازه (پوش)} یک بردار تصادفی گاوسی دوبعدی را مدل می‌کند. هنگامی که \lr{$X$} و \lr{$Y$} مستقل و \lr{$N(0,\, \sigma^{2})$} باشند، آنگاه \lr{$R = \sqrt{X^{2} + Y^{2}}$} از توزیع رایلی پیروی می‌کند.

\PSsubsection{زمینه مهندسی}

\begin{itemize}
\tightlist
\item
  پوش نویز باریک‌باند (هم‌فاز + مربعی)
\item
  اندازه ضریب کانال محوشدگی بی‌سیم (محوشدگی رایلی)
\item
  مدل‌سازی سرعت باد
\item
  خطای فاصله در مکان‌یابی دوبعدی (هنگامی که خطاها در \lr{$x$} و \lr{$y$} گاوسی هستند)
\end{itemize}

\begin{PSbox}[unbreakable]{def}{تعریف ریاضی}

\lr{$R \sim \mathrm{Rayleigh}(\sigma)$}، که در آن \lr{$\sigma$} پارامتر است (نه انحراف معیار \lr{$R$}).

\end{PSbox}

\PSsubsection{\lr{\mbox{PDF}}}

\PSdisplay{f_R(r) = \frac{r}{\sigma^2} e^{-r^2/(2\sigma^2)}, \quad r \geq 0}

\PSsubsection{\lr{\mbox{CDF}}}

\PSdisplay{F_R(r) = 1 - e^{-r^2/(2\sigma^2)}, \quad r \geq 0}

\PSsubsection{میانگین}

\PSdisplay{E[R] = \sigma\sqrt{\frac{\pi}{2}} \approx 1.2533\sigma}

\PSsubsection{واریانس}

\PSdisplay{\PSmtext{Var}(R) = \frac{4-\pi}{2}\sigma^2 \approx 0.4292\sigma^2}

\PSsubsection{رابطه با سایر توزیع‌ها}

\begin{itemize}
\tightlist
\item
  \lr{$R^{2} \sim \mathrm{Exponential}(2\sigma^{2})$} \lr{—} توان به‌صورت نمایی توزیع شده است
\item
  اگر \lr{$X,\, Y \sim N(0,\, \sigma^{2})$} مستقل باشند، آنگاه \lr{$\sqrt{X^{2}+Y^{2}} \sim \mathrm{Rayleigh}(\sigma)$}
\end{itemize}

\PSsubsection{زمان استفاده}

\lr{\mbox{\PSyes{}}} پوش/اندازه یک سیگنال گاوسی مختلط\PSbr{}\lr{\mbox{\PSyes{}}} مدل‌سازی کانال محوشدگی رایلی (بدون خط دید)\PSbr{}\lr{\mbox{\PSyes{}}} فاصله دوبعدی با خطاهای گاوسی در هر بُعد\PSbr{}\lr{\mbox{\PSyes{} RSS}} (جذر مجموع مربعات) دو مؤلفه گاوسی مستقل

\PSsubsection{زمان عدم استفاده}

\lr{\mbox{\PSno{}}} مؤلفه خط دید قوی موجود است (از رایسی استفاده کنید)\PSbr{}\lr{\mbox{\PSno{}}} اندازه/پوش را مدل نمی‌کنید\PSbr{}\lr{\mbox{\PSno{}}} مؤلفه‌ها گاوسی نیستند یا واریانس برابر ندارند

\PSsubsection{پیاده‌سازی در \lr{\mbox{MATLAB}}}

\tcbset{PScodemode/.style={unbreakable}}

\begin{Shaded}
\begin{Highlighting}[]
\CommentTok{\%\% Rayleigh Distribution: Wireless Fading Channel}
\VariableTok{sigma} \OperatorTok{=} \FloatTok{1}\OperatorTok{;}   \CommentTok{\% Rayleigh parameter}
\VariableTok{r} \OperatorTok{=} \VariableTok{linspace}\NormalTok{(}\FloatTok{0}\OperatorTok{,} \FloatTok{5}\OperatorTok{,} \FloatTok{500}\NormalTok{)}\OperatorTok{;}

\VariableTok{pdf\_ray} \OperatorTok{=} \VariableTok{raylpdf}\NormalTok{(}\VariableTok{r}\OperatorTok{,} \VariableTok{sigma}\NormalTok{)}\OperatorTok{;}
\VariableTok{cdf\_ray} \OperatorTok{=} \VariableTok{raylcdf}\NormalTok{(}\VariableTok{r}\OperatorTok{,} \VariableTok{sigma}\NormalTok{)}\OperatorTok{;}

\VariableTok{figure}\OperatorTok{;}
\VariableTok{subplot}\NormalTok{(}\FloatTok{2}\OperatorTok{,}\FloatTok{1}\OperatorTok{,}\FloatTok{1}\NormalTok{)}\OperatorTok{;}
\VariableTok{plot}\NormalTok{(}\VariableTok{r}\OperatorTok{,} \VariableTok{pdf\_ray}\OperatorTok{,} \SpecialStringTok{\textquotesingle{}b{-}\textquotesingle{}}\OperatorTok{,} \SpecialStringTok{\textquotesingle{}LineWidth\textquotesingle{}}\OperatorTok{,} \FloatTok{2}\NormalTok{)}\OperatorTok{;}
\VariableTok{xlabel}\NormalTok{(}\SpecialStringTok{\textquotesingle{}Channel Gain |h|\textquotesingle{}}\NormalTok{)}\OperatorTok{;} \VariableTok{ylabel}\NormalTok{(}\SpecialStringTok{\textquotesingle{}f\_R(r)\textquotesingle{}}\NormalTok{)}\OperatorTok{;}
\VariableTok{title}\NormalTok{(}\VariableTok{sprintf}\NormalTok{(}\SpecialStringTok{\textquotesingle{}Rayleigh PDF (σ = \%d)\textquotesingle{}}\OperatorTok{,} \VariableTok{sigma}\NormalTok{))}\OperatorTok{;}
\VariableTok{grid} \VariableTok{on}\OperatorTok{;}

\VariableTok{subplot}\NormalTok{(}\FloatTok{2}\OperatorTok{,}\FloatTok{1}\OperatorTok{,}\FloatTok{2}\NormalTok{)}\OperatorTok{;}
\VariableTok{plot}\NormalTok{(}\VariableTok{r}\OperatorTok{,} \VariableTok{cdf\_ray}\OperatorTok{,} \SpecialStringTok{\textquotesingle{}r{-}\textquotesingle{}}\OperatorTok{,} \SpecialStringTok{\textquotesingle{}LineWidth\textquotesingle{}}\OperatorTok{,} \FloatTok{2}\NormalTok{)}\OperatorTok{;}
\VariableTok{xlabel}\NormalTok{(}\SpecialStringTok{\textquotesingle{}r\textquotesingle{}}\NormalTok{)}\OperatorTok{;} \VariableTok{ylabel}\NormalTok{(}\SpecialStringTok{\textquotesingle{}P(R ≤ r)\textquotesingle{}}\NormalTok{)}\OperatorTok{;}
\VariableTok{title}\NormalTok{(}\SpecialStringTok{\textquotesingle{}Rayleigh CDF = Outage Probability\textquotesingle{}}\NormalTok{)}\OperatorTok{;}
\VariableTok{grid} \VariableTok{on}\OperatorTok{;}

\CommentTok{\% Outage probability: P(|h| \textless{} threshold)}
\VariableTok{threshold} \OperatorTok{=} \FloatTok{0.5}\OperatorTok{;}
\VariableTok{p\_outage} \OperatorTok{=} \VariableTok{raylcdf}\NormalTok{(}\VariableTok{threshold}\OperatorTok{,} \VariableTok{sigma}\NormalTok{)}\OperatorTok{;}
\VariableTok{fprintf}\NormalTok{(}\SpecialStringTok{\textquotesingle{}P(outage, |h| \textless{} \%.1f) = \%.4f\textbackslash{}n\textquotesingle{}}\OperatorTok{,} \VariableTok{threshold}\OperatorTok{,} \VariableTok{p\_outage}\NormalTok{)}\OperatorTok{;}
\end{Highlighting}
\end{Shaded}

\begin{PSbox}[unbreakable]{ex}{مثال شبیه‌سازی: محوشدگی رایلی}

\tcbset{PScodemode/.style={unbreakable}}

\begin{Shaded}
\begin{Highlighting}[]
\CommentTok{\%\% Generate Rayleigh Fading from Complex Gaussian}
\VariableTok{sigma} \OperatorTok{=} \FloatTok{1}\OperatorTok{;}
\VariableTok{N} \OperatorTok{=} \FloatTok{100000}\OperatorTok{;}

\CommentTok{\% Complex Gaussian channel coefficient}
\VariableTok{h} \OperatorTok{=} \VariableTok{sigma}\OperatorTok{/}\VariableTok{sqrt}\NormalTok{(}\FloatTok{2}\NormalTok{) }\OperatorTok{*}\NormalTok{ (}\VariableTok{randn}\NormalTok{(}\FloatTok{1}\OperatorTok{,}\VariableTok{N}\NormalTok{) }\OperatorTok{+} \FloatTok{1j}\OperatorTok{*}\VariableTok{randn}\NormalTok{(}\FloatTok{1}\OperatorTok{,}\VariableTok{N}\NormalTok{))}\OperatorTok{;}

\CommentTok{\% Magnitude is Rayleigh distributed}
\VariableTok{R} \OperatorTok{=} \VariableTok{abs}\NormalTok{(}\VariableTok{h}\NormalTok{)}\OperatorTok{;}

\VariableTok{figure}\OperatorTok{;}
\VariableTok{histogram}\NormalTok{(}\VariableTok{R}\OperatorTok{,} \FloatTok{100}\OperatorTok{,} \SpecialStringTok{\textquotesingle{}Normalization\textquotesingle{}}\OperatorTok{,} \SpecialStringTok{\textquotesingle{}pdf\textquotesingle{}}\NormalTok{)}\OperatorTok{;}
\VariableTok{hold} \VariableTok{on}\OperatorTok{;}
\VariableTok{r} \OperatorTok{=} \VariableTok{linspace}\NormalTok{(}\FloatTok{0}\OperatorTok{,} \FloatTok{5}\OperatorTok{,} \FloatTok{300}\NormalTok{)}\OperatorTok{;}
\VariableTok{plot}\NormalTok{(}\VariableTok{r}\OperatorTok{,} \VariableTok{raylpdf}\NormalTok{(}\VariableTok{r}\OperatorTok{,} \VariableTok{sigma}\OperatorTok{/}\VariableTok{sqrt}\NormalTok{(}\FloatTok{2}\NormalTok{))}\OperatorTok{,} \SpecialStringTok{\textquotesingle{}r{-}\textquotesingle{}}\OperatorTok{,} \SpecialStringTok{\textquotesingle{}LineWidth\textquotesingle{}}\OperatorTok{,} \FloatTok{2}\NormalTok{)}\OperatorTok{;}
\VariableTok{xlabel}\NormalTok{(}\SpecialStringTok{\textquotesingle{}|h|\textquotesingle{}}\NormalTok{)}\OperatorTok{;} \VariableTok{ylabel}\NormalTok{(}\SpecialStringTok{\textquotesingle{}PDF\textquotesingle{}}\NormalTok{)}\OperatorTok{;}
\VariableTok{title}\NormalTok{(}\SpecialStringTok{\textquotesingle{}Rayleigh Fading: Simulated vs Theory\textquotesingle{}}\NormalTok{)}\OperatorTok{;}
\VariableTok{legend}\NormalTok{(}\SpecialStringTok{\textquotesingle{}Simulation\textquotesingle{}}\OperatorTok{,} \SpecialStringTok{\textquotesingle{}Theory\textquotesingle{}}\NormalTok{)}\OperatorTok{;}
\end{Highlighting}
\end{Shaded}

\end{PSbox}

\PSsection{3.12}{توزیع کای‌دو}

\PSsubsection{تفسیر فیزیکی}

\textbf{جمع مربعات} متغیرهای تصادفی نرمال استاندارد مستقل را مدل می‌کند. به‌طور طبیعی در برآورد واریانس و آزمون نیکویی برازش ظاهر می‌شود.

\PSsubsection{زمینه مهندسی}

\begin{itemize}
\tightlist
\item
  توزیع واریانس نمونه (برای داده‌های گاوسی)
\item
  توان نویز گاوسی (جمع مؤلفه‌های مربع‌شده)
\item
  آماره آزمون نیکویی برازش
\item
  بازه‌های اطمینان برای واریانس
\item
  آشکارسازی انرژی در سنجش طیف
\end{itemize}

\begin{PSbox}[unbreakable]{def}{تعریف ریاضی}

اگر \lr{$Z_{1},\, Z_{2},\, \ldots,\, Z_{n}$} مستقل و \lr{$N(0,\,1)$} باشند، آنگاه:\PSdisplay{\chi^2 = Z_1^2 + Z_2^2 + \cdots + Z_n^2 \sim \chi^2(n)}

\lr{$n$} = درجات آزادی.

\end{PSbox}

\PSsubsection{\lr{\mbox{PDF}}}

\PSdisplay{f_X(x) = \frac{x^{n/2-1} e^{-x/2}}{2^{n/2}\Gamma(n/2)}, \quad x \geq 0}

\PSsubsection{\lr{\mbox{CDF}}}

برای \lr{$n$} کلی فرم بسته ندارد. از طریق تابع گامای ناقص محاسبه می‌شود.

\PSsubsection{میانگین}

\PSdisplay{E[X] = n}

\PSsubsection{واریانس}

\PSdisplay{\PSmtext{Var}(X) = 2n}

\PSsubsection{رابطه با سایر توزیع‌ها}

\begin{itemize}
\tightlist
\item
  \lr{Chi-square(n) $= \mathrm{Gamma}\left(\frac{n}{2},\, 2\right)$}
\item
  \lr{\mbox{Chi-square(1)}} = مربع \lr{$N(0,\,1)$}
\item
  \lr{Chi-square(2) $= \mathrm{Exponential}\left(\frac{1}{2}\right)$}
\item
  برای \lr{$n$} بزرگ: \lr{$\chi^{2}(n) \approx N(n,\, 2n)$} بر اساس \lr{\mbox{CLT}}
\end{itemize}

\PSsubsection{زمان استفاده}

\lr{\mbox{\PSyes{}}} جمع مربعات متغیرهای تصادفی گاوسی\PSbr{}\lr{\mbox{\PSyes{}}} آزمون واریانس و بازه‌های اطمینان\PSbr{}\lr{\mbox{\PSyes{}}} آزمون‌های نیکویی برازش\PSbr{}\lr{\mbox{\PSyes{}}} آشکارسازی انرژی در \lr{$N$} بُعد

\PSsubsection{زمان عدم استفاده}

\lr{\mbox{\PSno{}}} متغیرها گاوسی نیستند\PSbr{}\lr{\mbox{\PSno{}}} متغیرها مستقل نیستند\PSbr{}\lr{\mbox{\PSno{}}} متغیرها نرمال استاندارد نیستند (نیاز به مقیاس‌بندی دارند)

\PSsubsection{پیاده‌سازی در \lr{\mbox{MATLAB}}}

\tcbset{PScodemode/.style={unbreakable}}

\begin{Shaded}
\begin{Highlighting}[]
\CommentTok{\%\% Chi{-}Square Distribution: Noise Power}
\CommentTok{\% Noise power from N complex dimensions = sum of 2N squared Gaussians}

\VariableTok{n\_dof} \OperatorTok{=} \FloatTok{10}\OperatorTok{;}        \CommentTok{\% degrees of freedom}
\VariableTok{x} \OperatorTok{=} \VariableTok{linspace}\NormalTok{(}\FloatTok{0}\OperatorTok{,} \FloatTok{30}\OperatorTok{,} \FloatTok{500}\NormalTok{)}\OperatorTok{;}

\VariableTok{pdf\_chi2} \OperatorTok{=} \VariableTok{chi2pdf}\NormalTok{(}\VariableTok{x}\OperatorTok{,} \VariableTok{n\_dof}\NormalTok{)}\OperatorTok{;}
\VariableTok{cdf\_chi2} \OperatorTok{=} \VariableTok{chi2cdf}\NormalTok{(}\VariableTok{x}\OperatorTok{,} \VariableTok{n\_dof}\NormalTok{)}\OperatorTok{;}

\VariableTok{figure}\OperatorTok{;}
\VariableTok{plot}\NormalTok{(}\VariableTok{x}\OperatorTok{,} \VariableTok{pdf\_chi2}\OperatorTok{,} \SpecialStringTok{\textquotesingle{}b{-}\textquotesingle{}}\OperatorTok{,} \SpecialStringTok{\textquotesingle{}LineWidth\textquotesingle{}}\OperatorTok{,} \FloatTok{2}\NormalTok{)}\OperatorTok{;}
\VariableTok{xlabel}\NormalTok{(}\SpecialStringTok{\textquotesingle{}χ² Value (Noise Power)\textquotesingle{}}\NormalTok{)}\OperatorTok{;}
\VariableTok{ylabel}\NormalTok{(}\SpecialStringTok{\textquotesingle{}f\_X(x)\textquotesingle{}}\NormalTok{)}\OperatorTok{;}
\VariableTok{title}\NormalTok{(}\VariableTok{sprintf}\NormalTok{(}\SpecialStringTok{\textquotesingle{}Chi{-}Square PDF: \%d degrees of freedom\textquotesingle{}}\OperatorTok{,} \VariableTok{n\_dof}\NormalTok{))}\OperatorTok{;}
\VariableTok{grid} \VariableTok{on}\OperatorTok{;}

\CommentTok{\% Simulation: sum of squared standard normals}
\VariableTok{N} \OperatorTok{=} \FloatTok{100000}\OperatorTok{;}
\VariableTok{Z} \OperatorTok{=} \VariableTok{randn}\NormalTok{(}\VariableTok{n\_dof}\OperatorTok{,} \VariableTok{N}\NormalTok{)}\OperatorTok{;}
\VariableTok{X\_sim} \OperatorTok{=} \VariableTok{sum}\NormalTok{(}\VariableTok{Z}\OperatorTok{.\^{}}\FloatTok{2}\OperatorTok{,} \FloatTok{1}\NormalTok{)}\OperatorTok{;}

\VariableTok{fprintf}\NormalTok{(}\SpecialStringTok{\textquotesingle{}Mean: Theory=\%d, Sim=\%.2f\textbackslash{}n\textquotesingle{}}\OperatorTok{,} \VariableTok{n\_dof}\OperatorTok{,} \VariableTok{mean}\NormalTok{(}\VariableTok{X\_sim}\NormalTok{))}\OperatorTok{;}
\VariableTok{fprintf}\NormalTok{(}\SpecialStringTok{\textquotesingle{}Var:  Theory=\%d, Sim=\%.2f\textbackslash{}n\textquotesingle{}}\OperatorTok{,} \FloatTok{2}\OperatorTok{*}\VariableTok{n\_dof}\OperatorTok{,} \VariableTok{var}\NormalTok{(}\VariableTok{X\_sim}\NormalTok{))}\OperatorTok{;}
\end{Highlighting}
\end{Shaded}

\PSsection{3.13}{راهنمای انتخاب توزیع}

\PSsubsection{درخت تصمیم: توزیع‌های گسسته}

\PSfigure{m03-tree-discrete}

\PSsubsection{درخت تصمیم: توزیع‌های پیوسته}

\PSfigure{m03-tree-continuous}

\PSsubsection{الگوهای تشخیص سریع}

{\def\LTcaptype{none} % do not increment counter
\begin{longtable}[c]{@{}cc@{}}
\toprule\noalign{}
\rowcolor{PSnavy}\PSth{اگر ببینید\lr{…}} & \PSth{فکر کنید\lr{…}} \\
\midrule\noalign{}
\endhead
\bottomrule\noalign{}
\endlastfoot
خروجی دودویی & برنولی \\
«\lr{$n$} آزمایش، \lr{$k$} موفقیت» & دوجمله‌ای \\
«چند تا تا نخستین\lr{…}» & هندسی \\
«رویداد در هر بازه» & پواسون \\
«شانس برابر در هر جای \lr{$[a,\,b]$}» & یکنواخت \\
«زمان تا رویداد بعدی» & نمایی \\
«جمع بسیاری از اثرها» یا «نویز» & گاوسی \\
«کل زمان انتظار برای \lr{$k$} رویداد» & گاما \\
«پوش سیگنال» یا «محوشدگی» & رایلی \\
«جمع مربعات نرمال‌ها» & کای‌دو \\
\end{longtable}
}

\PSsection{3.14}{جدول خلاصه}

{\def\LTcaptype{none} % do not increment counter
\begin{longtable}[c]{@{}
  >{\centering\arraybackslash}p{(0.965\linewidth - 10\tabcolsep) * \real{0.2407}}
  >{\centering\arraybackslash}p{(0.965\linewidth - 10\tabcolsep) * \real{0.1111}}
  >{\centering\arraybackslash}p{(0.965\linewidth - 10\tabcolsep) * \real{0.2037}}
  >{\centering\arraybackslash}p{(0.965\linewidth - 10\tabcolsep) * \real{0.1111}}
  >{\centering\arraybackslash}p{(0.965\linewidth - 10\tabcolsep) * \real{0.1852}}
  >{\centering\arraybackslash}p{(0.965\linewidth - 10\tabcolsep) * \real{0.1481}}@{}}
\toprule\noalign{}
\rowcolor{PSnavy}\PSth{توزیع} & \PSth{نوع} & \PSth{پارامترها} & \PSth{میانگین} & \PSth{واریانس} & \PSth{\lr{\mbox{MATLAB}}} \\
\midrule\noalign{}
\endhead
\bottomrule\noalign{}
\endlastfoot
برنولی & گسسته & \lr{$p$} & \lr{$p$} & \lr{$p(1-p)$} & \lr{\mbox{\PScode{binornd(1,p)}}} \\
دوجمله‌ای & گسسته & \lr{$n,\, p$} & \lr{$np$} & \lr{$np(1-p)$} & \lr{\mbox{\PScode{binornd(n,p)}}} \\
هندسی & گسسته & \lr{$p$} & \lr{$\frac{1}{p}$} & \lr{$\frac{1-p}{p^{2}}$} & \lr{\mbox{\PScode{geornd(p)+1}}} \\
دوجمله‌ای منفی & گسسته & \lr{$r,\, p$} & \lr{$\frac{r}{p}$} & \lr{$\frac{r(1-p)}{p^{2}}$} & \lr{\mbox{\PScode{nbinrnd(r,p)+r}}} \\
پواسون & گسسته & \lr{$\lambda$} & \lr{$\lambda$} & \lr{$\lambda$} & \lr{\mbox{\PScode{poissrnd(λ)}}} \\
یکنواخت & پیوسته & \lr{$a,\, b$} & \lr{$\frac{a+b}{2}$} & \lr{$\frac{(b-a)^{2}}{12}$} & \lr{\mbox{\PScode{unifrnd(a,b)}}} \\
نمایی & پیوسته & \lr{$\lambda$} (یا \lr{$\beta = \frac{1}{\lambda}$}) & \lr{$\frac{1}{\lambda}$} & \lr{$\frac{1}{\lambda^{2}}$} & \lr{\mbox{\PScode{exprnd(β)}}} \\
گاوسی & پیوسته & \lr{$\mu,\, \sigma^{2}$} & \lr{$\mu$} & \lr{$\sigma^{2}$} & \lr{\mbox{\PScode{normrnd(μ,σ)}}} \\
گاما & پیوسته & \lr{$\alpha,\, \beta$} & \lr{$\alpha \beta$} & \lr{$\alpha \beta^{2}$} & \lr{\mbox{\PScode{gamrnd(α,β)}}} \\
رایلی & پیوسته & \lr{$\sigma$} & \lr{$\sigma \sqrt{\frac{\pi}{2}}$} & \lr{$\frac{(4-\pi)\sigma^{2}}{2}$} & \lr{\mbox{\PScode{raylrnd(σ)}}} \\
کای‌دو & پیوسته & \lr{$n$} & \lr{$n$} & \lr{$2n$} & \lr{\mbox{\PScode{chi2rnd(n)}}} \\
\end{longtable}
}

\PSsection{3.15}{مسایل تمرینی}

\begin{PSbox}[unbreakable]{prob}{مسیله 1: تشخیص توزیع}

توزیع مناسب را برای هر سناریو تشخیص دهید:

\begin{enumerate}
\def\labelenumi{(\PSalph{enumi})}
\tightlist
\item
  تعداد تماس‌های قطع‌شده در یک پنجره 10 دقیقه‌ای (میانگین: 2 در هر 10 دقیقه)
\item
  ولتاژ نویز حرارتی در دو سر یک مقاومت
\item
  زمان تا از کار افتادن یک سرور (نرخ خرابی ثابت)
\item
  تعداد مدارهای مجتمع معیوب از میان یک دسته 100 تایی (نرخ نقص 3\%)
\item
  پوش سیگنال دریافتی از مسیر چندگانه بدون \lr{\mbox{LOS}}
\end{enumerate}

\PSsolution{} (الف) \lr{$\mathrm{Poisson}(\lambda = 2)$}، (ب) \lr{$\mathrm{Gaussian}(0,\, \sigma^{2})$}، (ج) \lr{$\mathrm{Exponential}(\lambda)$}، (د) \lr{$\mathrm{Binomial}(100,\, 0.03)$}، (ه) \lr{$\mathrm{Rayleigh}(\sigma)$}

\end{PSbox}

\begin{PSbox}[unbreakable]{prob}{مسیله 2: نمایی در برابر پواسون}

بسته‌ها به‌صورت یک فرایند پواسون با نرخ \lr{$\lambda = 5$} بسته بر ثانیه به یک مسیریاب می‌رسند.

\begin{enumerate}
\def\labelenumi{(\PSalph{enumi})}
\tightlist
\item
  \lr{$P(\text{\rl{بیش از 8 بسته در 1 ثانیه}})$} چقدر است؟
\item
  \lr{$P(\text{\rl{هیچ بسته‌ای در 0.5 ثانیه بعدی نرسد}})$} چقدر است؟
\item
  چه توزیعی زمان بین‌ورودی را توصیف می‌کند؟
\item
  میانگین زمان بین‌ورودی چقدر است؟
\end{enumerate}

\PSsolution{} 

\begin{enumerate}
\def\labelenumi{(\PSalph{enumi})}
\tightlist
\item
  \lr{$X \sim \mathrm{Poisson}(5)$: $P(X > 8) = 1 - \texttt{poisscdf}(8,\, 5) = 0.0681$}
\item
  \lr{$T \sim \mathrm{Exp}(5)$: $P(T > 0.5) = e^{-5 \times 0.5} = e^{-2.5} = 0.0821$}
\item
  نمایی با نرخ \lr{$\lambda = 5$}
\item
  میانگین \lr{$= \frac{1}{\lambda} = 0.2$} ثانیه
\end{enumerate}

\end{PSbox}

\begin{PSbox}[unbreakable]{prob}{مسیله 3: احتمال گاوسی}

یک مقاومت مقدار نامی \lr{\mbox{1000\ensuremath{\Omega}}} با تغییرات ساخت مدل‌شده به‌صورت \lr{$N(1000,\, 25)$} دارد \lr{$(\sigma = 5\,\Omega)$}.

\begin{enumerate}
\def\labelenumi{(\PSalph{enumi})}
\tightlist
\item
  چه کسری از مقاومت‌ها در محدوده \lr{$\pm 10\,\Omega$} حول مقدار نامی هستند (مشخصه: \lr{$990-1010\,\Omega$})؟
\item
  چه کسری در آزمون رواداری \lr{$\pm 1\%$} رد می‌شوند؟
\item
  کدام بازه رواداری شامل \PSnum{99.7\%} مقاومت‌ها است؟
\end{enumerate}

\PSsolution{} 

\begin{enumerate}
\def\labelenumi{(\PSalph{enumi})}
\tightlist
\item
  \lr{$P(990 \le X \le 1010) = P(-2 \le Z \le 2) = 0.9545 \to 95.45\%$}
\item
  \lr{\mbox{$\pm 1\% = \pm 10\,\Omega$ \ensuremath{\to}}} همانند \lr{\mbox{(a)}}، نرخ رد \lr{$= 1 - 0.9545 = 4.55\%$}
\item
  \lr{\mbox{$\pm 3\sigma = \pm 15\,\Omega$ \ensuremath{\to}}} محدوده \lr{$[985,\, 1015]\Omega$}
\end{enumerate}

\emph{ماژول بعدی: {ماژول 4 \lr{—} امید ریاضی و گشتاورها}}

\end{PSbox}
